\documentclass[12pt,a4paper]{article}

\usepackage{iftex}
\ifPDFTeX
  \usepackage[utf8]{inputenc}
  \usepackage[T5]{fontenc}
  \usepackage[vietnamese]{babel}
  \usepackage{lmodern}
\else
  % phương án dự phòng khi biên dịch bằng XeLaTeX/LuaLaTeX
  \usepackage{fontspec}
  \setmainfont{Latin Modern Roman}\setsansfont{Latin Modern Sans}\setmonofont{Latin Modern Mono}
  \renewcommand{\contentsname}{Mục lục}\renewcommand{\figurename}{Hình}\renewcommand{\tablename}{Bảng}
\fi
\usepackage[margin=2.2cm]{geometry}
\usepackage{amssymb}
\usepackage{graphicx}
\usepackage{float}
\usepackage{tikz}
\usetikzlibrary{calc,shapes.geometric,arrows.meta,positioning}
\usepackage{booktabs}
\usepackage{longtable}
\usepackage{tabularx}
\usepackage{array}
\usepackage[table]{xcolor}
\usepackage{listings}
\usepackage{enumitem}
\usepackage{fancyhdr}
\usepackage{titlesec}
\usepackage[hidelinks]{hyperref}

\definecolor{bkblue}{RGB}{3,43,145}
\definecolor{codebg}{RGB}{246,248,250}
\newcommand{\todo}[1]{\textcolor{red}{\textbf{[#1]}}}
\newcolumntype{L}[1]{>{\raggedright\arraybackslash}p{#1}}

\lstset{
  basicstyle=\ttfamily\footnotesize, backgroundcolor=\color{codebg},
  frame=single, rulecolor=\color{black!20},
  breaklines=true, showstringspaces=false, columns=fullflexible
}

\titleformat{\section}{\Large\bfseries\color{bkblue}}{\thesection.}{0.6em}{}
\titleformat{\subsection}{\large\bfseries}{\thesubsection.}{0.6em}{}

\pagestyle{fancy}
\fancyhf{}
\lhead{Kiểm tra phần mềm -- Nhóm HEHE2027}
\rhead{Bài tập lớn 2: Black-box testing}
\cfoot{\thepage}
\setlength{\headheight}{15pt}

\renewcommand{\arraystretch}{1.2}
\sloppy
\setlength{\emergencystretch}{3em}

% Macro khung cho các feature do thành viên khác hoàn thiện
\newcommand{\featureskeleton}[3]{%
  \subsection{Feature #1: #2}
  \textbf{Mô tả sơ bộ:} #3 \\
  \todo{Thành viên phụ trách bổ sung đầy đủ theo đúng cấu trúc của Feature 001: mô tả và hướng dẫn sử dụng; use case + activity diagram;
  bảng biến đầu vào; ECP; BVA; decision table; use-case testing; danh sách test case \texttt{TC-#1-YYY}. Phải áp dụng đủ cả bốn kỹ thuật và xét mọi biến đầu vào.}
}

\begin{document}

% ================= COVER =================
\begin{titlepage}
\begin{tikzpicture}[overlay,remember picture]
\draw [line width=2.25pt]
    ($ (current page.north west) + (2.0cm,-1.5cm) $)
    rectangle
    ($ (current page.south east) + (-1.5cm,2cm) $);
\draw [line width=0.5pt]
    ($ (current page.north west) + (2.2cm,-1.7cm) $)
    rectangle
    ($ (current page.south east) + (-1.7cm,2.2cm) $);
\end{tikzpicture}
\begin{center}
\textbf{ĐẠI HỌC QUỐC GIA THÀNH PHỐ HỒ CHÍ MINH} \\
\textbf{TRƯỜNG ĐẠI HỌC BÁCH KHOA}\\
\textbf{KHOA KHOA HỌC VÀ KỸ THUẬT MÁY TÍNH}

\vspace{1cm}
\includegraphics[width=0.32\textwidth]{img/logo/logo.png}
\vspace{1cm}

\begin{tabular}{c}
    \textbf{\Large KIỂM TRA PHẦN MỀM (SOFTWARE TESTING)}\\[0.4cm]
    \textbf{\large BÀI TẬP LỚN 2: BLACK-BOX TESTING}\\[0.2cm]
    \large Kiểm thử hệ thống quản lý học tập \textit{Moodle LMS}\\[0.4cm]
    \textbf{Học kỳ 1, năm học 2026--2027 -- Lớp L01 -- Nhóm HEHE2027}\\
    \hline
\end{tabular}

\vspace{1cm}
\begin{tabular}{p{3cm} p{6cm} l}
    \textbf{GVHD:} & Trần Thị Ngọc Trâm & \\[0.2cm]
    \textbf{Sinh viên:} & Nguyễn Đức Phương & 2312749 \\
                        & Lê Sĩ Nhật Quang & 2312791 \\
                        & Huỳnh Đức Tài & 2212971 \\
                        & Đào Xuân Sơn & 2312954 \\
                        & Sơn Văn Phát & 2212535 \\
\end{tabular}

\vfill
\fontsize{14pt}{0pt}\selectfont TP. HỒ CHÍ MINH, THÁNG 11/2026
\end{center}
\end{titlepage}

\tableofcontents
\newpage

% ================= 1 =================
\section{Giới thiệu}
Nhóm áp dụng các kỹ thuật kiểm thử hộp đen để thiết kế test case cho \textbf{Moodle LMS} (bản Docker theo
hướng dẫn của đề bài), gồm: phân tích giá trị biên (\textit{Boundary Value Analysis -- BVA}),
phân lớp tương đương (\textit{Equivalence Class Partitioning -- ECP}), bảng quyết định
(\textit{Decision Table}, phần cộng điểm) và kiểm thử theo use case (\textit{Use-case testing}).
Các test case được ghi trong file \texttt{.xlsx} theo mẫu của đề bài và được tự động hóa bằng \textit{Katalon Recorder}.

Mỗi thành viên chọn một feature riêng, không trùng nhau và không thuộc nhóm bị cấm (đăng nhập, đăng xuất, đổi mật khẩu
user mặc định). Mã test case có dạng \texttt{TC-XXX-YYY} với XXX là số thứ tự feature và YYY là số thứ tự test case.

\textbf{Tài liệu nộp kèm:} (1) báo cáo này (\texttt{.pdf}); (2) file \texttt{HEHE2027\_L01\_Proj2\_TestCases.xlsx};
(3) file nén \texttt{.zip} chứa các tệp \texttt{.krecorder}.

\subsection*{Quy định bổ sung của đề cần tuân thủ}
\begin{enumerate}[noitemsep]
  \item Mọi kỹ thuật phải được áp dụng cho \textbf{từng} feature. Nếu một kỹ thuật không áp dụng được về mặt logic
        (ví dụ bảng quyết định cho máy tính BMI, vì không có tổ hợp đầu vào nào cho kết quả khác nhau) thì không chọn feature đó.
  \item Mỗi sinh viên phải thể hiện hiểu biết về \textbf{tất cả} kỹ thuật bắt buộc, nên feature phải được chọn cẩn thận
        ngay từ đầu.
  \item \textbf{Phải chạy test suite Katalon Recorder} (ít nhất trên một trình duyệt, ví dụ Chrome) \textbf{trước khi tải về};
        log chạy được tải kèm test suite.
  \item File \texttt{.xlsx} dành cho \textbf{kiểm thử thủ công}: người kiểm thử làm theo cột \textit{Step name} và
        \textit{Step description} rồi đối chiếu với \textit{Expected Result}; đây không phải script Katalon.
        Cột \textit{Precondition} nêu rõ điều kiện trước khi test (ví dụ: đã đăng nhập bằng user nào, ngôn ngữ trình duyệt là English).
  \item \textbf{Không đổi mật khẩu của các user mặc định} trong Moodle demo.
\end{enumerate}

% ================= 2 =================
\section{Phân công công việc}

{\small
\begin{longtable}{|c|L{3.2cm}|L{3.8cm}|L{4.8cm}|c|}
\hline
\rowcolor{gray!15}
\textbf{STT} & \textbf{Thành viên} & \textbf{Feature phụ trách} & \textbf{Công việc} & \textbf{Đóng góp} \\ \hline
\endhead
1 & Nguyễn Đức Phương \newline 2312749 & 001 -- Add a new user &
  Phân tích, thiết kế test case, viết script Katalon cho feature của mình; & 20\% \\ \hline
2 & Lê Sĩ Nhật Quang \newline 2312791 & 002 -- Add a new course & Như trên & 20\% \\ \hline
3 & Huỳnh Đức Tài \newline 2212971 & 003 -- Add a quiz & Như trên & 20\% \\ \hline
4 & Đào Xuân Sơn \newline 2312954 & 004 -- Add new category & Như trên & 20\% \\ \hline
5 & Sơn Văn Phát \newline 2212535 & 005 -- Enrol users & Như trên & 20\% \\ \hline
\end{longtable}
}

% ================= 3 =================
\section{Đặc tả yêu cầu: các feature được chọn}
Đối tượng kiểm thử: \textbf{Moodle LMS}, thao tác bằng tài khoản quản trị viên. Các feature được chọn đều có nhiều
trường nhập liệu có ràng buộc (để áp dụng ECP, BVA), có các điều kiện kết hợp (decision table) và có luồng nghiệp vụ
có nhánh (use-case testing).

\textbf{Tiêu chí chọn feature:} feature phải (1) có nhiều biến đầu vào có ràng buộc (BVA, ECP);
(2) có tổ hợp điều kiện cho ra kết quả khác nhau (bảng quyết định); (3) có luồng nghiệp vụ có nhánh (use-case testing);
(4) không thuộc nhóm bị cấm và không làm đổi mật khẩu user mặc định.

\begin{table}[H]
\centering
\small
\begin{tabular}{|c|L{3.8cm}|L{6.8cm}|L{2.4cm}|}
\hline
\rowcolor{gray!15}\textbf{No.} & \textbf{Feature} & \textbf{Mô tả ngắn} & \textbf{Tester} \\ \hline
001 & Add a new user & Quản trị viên tạo tài khoản người dùng mới. & Phương \\ \hline
002 & Add a new course & Quản trị viên tạo khóa học mới (full name, short name, category, visibility, start/end date, ID number). & Quang \\ \hline
003 & --- & --- & Tài \\ \hline
004 & Add new category & Quản trị viên tạo danh mục khóa học mới. & Sơn \\ \hline
005 & --- & --- & Phát \\ \hline
\end{tabular}
\caption{Danh sách feature}
\end{table}

\begin{table}[H]
\centering\small
\begin{tabular}{|c|L{3.2cm}|L{8.6cm}|}
\hline
\rowcolor{gray!15}\textbf{No.} & \textbf{Feature} & \textbf{Tổ hợp điều kiện dự kiến cho decision table} \\ \hline
001 & Add a new user & Generate password, password hợp lệ, username/email không trùng, Force password change, Suspended. \\ \hline
002 & Add a new course & Full name/short name hợp lệ, short name không trùng, đã chọn category, bật end date, end date $>$ start date. \\ \hline
003 & --- & --- \\ \hline
004 & Add new category & Parent category đã chọn, Name hợp lệ, Category ID number (rỗng/duy nhất/trùng), Parent là Top hay danh mục khác. \\ \hline
005 & --- & --- \\ \hline
\end{tabular}
\caption{Kiểm tra mỗi feature có áp dụng được bảng quyết định}
\end{table}

% ================= 4 =================
\section{Môi trường kiểm thử}
\begin{itemize}[noitemsep]
  \item Moodle chạy bằng Docker theo đề bài; địa chỉ truy cập \texttt{http://127.0.0.1:9501}, tài khoản quản trị
        có sẵn \texttt{admin}/\texttt{admin} (sửa nếu máy nhóm dùng cổng khác).
  \item Trình duyệt Chrome (hoặc Firefox) có cài \textit{Katalon Recorder}.
  \item Cấu hình demo (\texttt{assets/config.php}) đặt \texttt{passwordpolicy = 0}: Moodle không bắt mật khẩu tối thiểu 8 ký tự hay
        chữ hoa/số/ký tự đặc biệt; mật khẩu chỉ bị giới hạn tối đa 128 ký tự. Demo cũng đặt \texttt{noemailever = 1} (không gửi email thật).
  \item Đăng nhập chỉ được dùng như \textit{bước chuẩn bị} của test, không phải đối tượng kiểm thử.
  \item Ngôn ngữ giao diện Moodle và trình duyệt là English (tên menu và thông báo lỗi trong test case theo tiếng Anh).
  \item \textbf{Không thay đổi mật khẩu của các user mặc định} trong Moodle demo.
\end{itemize}

% ================= 5 =================
\section{Feature 001: Add a new user}

\subsection{Mô tả và hướng dẫn sử dụng}
Quản trị viên tạo tài khoản mới tại \textbf{Site administration $>$ Users $>$ Accounts $>$ Add a new user}
(URL: \texttt{<BASE\_URL>/user/editadvanced.php?id=-1}).

\textbf{Các bước sử dụng:}
\begin{enumerate}[noitemsep]
  \item Đăng nhập bằng tài khoản admin.
  \item Mở trang Add a new user.
  \item Nhập các trường bắt buộc (username, password hoặc chọn \textit{Generate password and notify user},
        first name, last name, email) và các trường tùy chọn (city, country, \dots).
  \item Bấm \textbf{Create user} (hoặc \textbf{Cancel} để hủy).
  \item Hệ thống kiểm tra dữ liệu; nếu hợp lệ thì tạo user và chuyển về \textit{Browse list of users}, nếu không thì
        hiển thị lỗi tại trường tương ứng và giữ nguyên form.
\end{enumerate}

\subsection{Use case và activity diagram}
\textbf{Use case:}
\begin{table}[H]
\centering\small
\begin{tabular}{|L{3.6cm}|L{11.4cm}|}
\hline
Tên & UC-001: Tạo người dùng mới \\ \hline
Tác nhân & Quản trị viên (Admin) \\ \hline
Tiền điều kiện & Admin đã đăng nhập; có quyền tạo người dùng. \\ \hline
Hậu điều kiện (thành công) & Tài khoản mới được lưu và hiển thị trong Browse list of users. \\ \hline
Luồng chính & Mở form $\rightarrow$ nhập dữ liệu hợp lệ $\rightarrow$ Create user $\rightarrow$ hệ thống tạo user và chuyển về danh sách. \\ \hline
Luồng thay thế & A1: thiếu hoặc sai dữ liệu $\rightarrow$ báo lỗi, nhập lại. A2: username hoặc email trùng $\rightarrow$ báo lỗi, đổi dữ liệu.
A3: chọn Generate password thay cho nhập password. A4: bấm Cancel, không tạo user. \\ \hline
\end{tabular}
\caption{Đặc tả use case UC-001}
\end{table}

\begin{figure}[H]
\centering
\begin{tikzpicture}[>=Stealth, node distance=0.9cm and 1.4cm,
  act/.style={rectangle, rounded corners=4pt, draw, minimum width=3.6cm, minimum height=0.8cm, align=center, font=\small, fill=blue!5},
  dec/.style={diamond, aspect=2.2, draw, align=center, font=\scriptsize, inner sep=1pt, fill=yellow!15},
  st/.style={circle, fill=black, minimum size=0.35cm},
  en/.style={circle, draw, double, minimum size=0.35cm}]
\node[st] (s) {};
\node[act, below=0.6cm of s] (a1) {Mở trang Add a new user};
\node[act, below=0.6cm of a1] (a2) {Nhập dữ liệu vào form};
\node[dec, below=0.7cm of a2] (d0) {Bấm\\Cancel?};
\node[act, below=0.8cm of d0] (a3) {Bấm Create user};
\node[dec, below=0.7cm of a3] (d1) {Dữ liệu\\hợp lệ?};
\node[dec, below=0.8cm of d1] (d2) {Username/email\\không trùng?};
\node[act, below=0.8cm of d2] (a4) {Tạo user, chuyển về\\Browse list of users};
\node[en, below=0.6cm of a4] (e) {};
\node[act, right=2.6cm of d0] (c1) {Quay lại danh sách,\\không tạo user};
\node[act, right=2.6cm of d1] (er1) {Hiển thị lỗi\\tại trường};
\node[act, right=2.6cm of d2] (er2) {Hiển thị lỗi\\trùng dữ liệu};
\draw[->] (s)--(a1); \draw[->] (a1)--(a2); \draw[->] (a2)--(d0);
\draw[->] (d0)--node[right,font=\scriptsize]{Không}(a3);
\draw[->] (d0)--node[above,font=\scriptsize]{Có}(c1);
\draw[->] (a3)--(d1);
\draw[->] (d1)--node[right,font=\scriptsize]{Có}(d2);
\draw[->] (d1)--node[above,font=\scriptsize]{Không}(er1);
\draw[->] (d2)--node[right,font=\scriptsize]{Có}(a4);
\draw[->] (d2)--node[above,font=\scriptsize]{Không}(er2);
\draw[->] (a4)--(e);
\draw[->] (er1.east)--++(1.0,0)|-(a2.east);
\draw[->] (er2.east)--++(1.6,0)|-(a2.east);
\draw[->] (c1.east)--++(2.4,0)|-(e.east);
\end{tikzpicture}
\caption{Activity diagram của UC-001 (các nhánh báo lỗi quay lại bước nhập dữ liệu; nhánh Cancel kết thúc)}
\end{figure}


\subsection{Các biến đầu vào}
\begin{table}[H]
\centering\small
\begin{tabular}{|L{2.6cm}|L{3.0cm}|L{4.2cm}|L{4.2cm}|}
\hline
\rowcolor{gray!15}\textbf{Biến} & \textbf{Bắt buộc?} & \textbf{Ràng buộc} & \textbf{Ghi chú} \\ \hline
Username & Có & 1--100 ký tự; chữ thường, số, \texttt{\_ - . @}; không trùng & Không phân biệt hoa/thường: chữ hoa bị từ chối \\ \hline
New password & Có (trừ khi chọn Generate password) & 1--128 ký tự; demo tắt password policy nên không yêu cầu độ phức tạp & Theo cấu hình \texttt{passwordpolicy = 0} của demo \\ \hline
First name & Có & 1--100 ký tự & \\ \hline
Last name & Có & 1--100 ký tự & \\ \hline
Email address & Có & Đúng định dạng, tối đa 100 ký tự, không trùng & \\ \hline
City/town & Không & 0--120 ký tự & \\ \hline
Country, Timezone, Description,\dots & Không & Chọn từ danh sách / văn bản tự do & Không có giá trị không hợp lệ qua giao diện; mỗi trường được điền trong TC-001-001 \\ \hline
Generate password and notify user & Không (checkbox) & Chọn/không chọn & Dùng trong decision table \\ \hline
Force password change & Không (checkbox) & Chọn/không chọn & Dùng trong decision table \\ \hline
Suspended account & Không (checkbox) & Chọn/không chọn & Dùng trong decision table \\ \hline
\end{tabular}
\caption{Các biến đầu vào của Feature 001}
\end{table}

\subsection{Kỹ thuật 1: Phân lớp tương đương (ECP)}
\textbf{Cách áp dụng:} với mỗi biến, chia miền giá trị thành lớp hợp lệ (V) và các lớp không hợp lệ (I).
Một test case ``tất cả hợp lệ'' (TC-001-001) phủ mọi lớp hợp lệ; mỗi lớp không hợp lệ được kiểm thử trong một test case riêng,
các trường khác giữ giá trị hợp lệ để cô lập nguyên nhân lỗi.

\begin{table}[H]
\centering\footnotesize
\begin{tabular}{|l|L{5.0cm}|L{7.4cm}|}
\hline
\rowcolor{gray!15}\textbf{Biến} & \textbf{Lớp hợp lệ} & \textbf{Lớp không hợp lệ (test case)} \\ \hline
Username & V1: 1--100 ký tự, chữ thường/số/ký tự cho phép, chưa tồn tại &
 I1: rỗng (002); I2: $>$100 ký tự (003); I3: có chữ hoa (004); I4: có khoảng trắng/ký tự cấm (005); I5: đã tồn tại (006) \\ \hline
Password & V2: 1--128 ký tự (demo không có password policy) & I6: rỗng, không chọn Generate password (007); I7: $>$128 ký tự (008) \\ \hline
First name & V3: 1--100 ký tự & I8: rỗng (009); I9: $>$100 (010) \\ \hline
Last name & V4: 1--100 ký tự & I10: rỗng (011); I11: $>$100 (012) \\ \hline
Email & V5: định dạng đúng, $\leq$100 ký tự, chưa tồn tại & I12: rỗng (013); I13: sai định dạng (014); I14: đã tồn tại (015); I15: $>$100 (016) \\ \hline
City & V6: 0--120 ký tự & I16: $>$120 (017) \\ \hline
\end{tabular}
\caption{Lớp tương đương và test case tương ứng (ID rút gọn: 002 nghĩa là TC-001-002)}
\end{table}

\subsection{Kỹ thuật 2: Phân tích giá trị biên (BVA)}
\textbf{Cách áp dụng:} với mỗi biến có miền độ dài liên tục, chọn các giá trị tại và quanh biên: min$-1$, min, min$+1$,
max$-1$, max, max$+1$. Các giá trị \emph{không hợp lệ} (min$-1$, max$+1$) trùng với test case ECP nên được tham chiếu lại,
không tạo thêm; các giá trị hợp lệ tạo test case mới (018--040). Password không có độ dài tối thiểu do demo tắt password policy,
nên min là 1 ký tự (rỗng bị báo ``Required'').

\begin{table}[H]
\centering\footnotesize
\begin{tabular}{|l|c|c|c|c|c|c|}
\hline
\rowcolor{gray!15}\textbf{Biến (độ dài)} & \textbf{min$-1$} & \textbf{min} & \textbf{min$+1$} & \textbf{max$-1$} & \textbf{max} & \textbf{max$+1$} \\ \hline
Username (1--100) & 0 (002) & 1 (018) & 2 (019) & 99 (020) & 100 (021) & 101 (003) \\ \hline
Password (1--128) & 0 (007) & 1 (022) & 2 (023) & 127 (024) & 128 (025) & 129 (008) \\ \hline
First name (1--100) & 0 (009) & 1 (026) & 2 (027) & 99 (028) & 100 (029) & 101 (010) \\ \hline
Last name (1--100) & 0 (011) & 1 (030) & 2 (031) & 99 (032) & 100 (033) & 101 (012) \\ \hline
Email (tổng độ dài) & rỗng (013) & a@b.co (034) & -- & 99 (035) & 100 (036) & 101 (016) \\ \hline
City (0--120) & -- & 0 (037) & 1 (038) & 119 (039) & 120 (040) & 121 (017) \\ \hline
\end{tabular}
\caption{Giá trị biên và test case (số trong ngoặc là TC-001-YYY)}
\end{table}

\subsection{Kỹ thuật 3: Bảng quyết định (Decision table)}
\textbf{Điều kiện:} C1 -- chọn Generate password; C2 -- đã nhập password (chỉ xét khi C1 = N; demo không có password policy);
C3 -- username và email hợp lệ, không trùng; C4 -- chọn Force password change; C5 -- chọn Suspended.
\textbf{Hành động:} A1 -- tạo user; A2 -- hiển thị lỗi; A3 -- đặt cờ force password change; A4 -- đặt trạng thái suspended.
Quy tắc R2 lấy từ mã nguồn: Moodle không cho chọn đồng thời Generate password và Suspended.

\begin{table}[H]
\centering\footnotesize
\begin{tabular}{|l|c|c|c|c|c|c|c|c|}
\hline
\rowcolor{gray!15}\textbf{Luật} & R1 & R2 & R3 & R4 & R5 & R6 & R7 & R8 \\ \hline
C1 Generate password & Y & Y & N & N & N & N & N & Y \\ \hline
C2 Đã nhập password & -- & -- & Y & Y & Y & N & Y & -- \\ \hline
C3 Username/email OK & Y & Y & Y & Y & Y & Y & N & N \\ \hline
C4 Force password change & N & -- & N & Y & N & -- & -- & -- \\ \hline
C5 Suspended & N & Y & N & N & Y & -- & -- & -- \\ \hline
A1 Tạo user & X & & X & X & X & & & \\ \hline
A2 Báo lỗi & & X & & & & X & X & X \\ \hline
A3 Force change & & & & X & & & & \\ \hline
A4 Suspended & & & & & X & & & \\ \hline
\textbf{Test case} & 041 & 042 & 001 & 043 & 044 & 007 & 006 & 045 \\ \hline
\end{tabular}
\caption{Decision table của Feature 001 (R3, R6, R7 tái sử dụng test case ECP)}
\end{table}
\textit{Ghi chú:} email thông báo mật khẩu (R1) không kiểm tra được vì demo đặt \texttt{noemailever = 1}; test chỉ kiểm tra
user được tạo. Tổ hợp Generate password cùng Force password change chưa được khảo sát nên không đưa vào bảng.

\subsection{Kỹ thuật 4: Use-case testing}
Từ activity diagram, các đường đi (scenario) được liệt kê như sau:

\begin{table}[H]
\centering\footnotesize
\begin{tabular}{|c|L{7.4cm}|L{3.2cm}|c|}
\hline
\rowcolor{gray!15}\textbf{Path} & \textbf{Đường đi trong diagram} & \textbf{Loại} & \textbf{Test case} \\ \hline
P1 & Mở form $\rightarrow$ nhập hợp lệ $\rightarrow$ Create user $\rightarrow$ tạo thành công & Luồng chính & 046 \\ \hline
P2 & Gửi form thiếu dữ liệu $\rightarrow$ báo lỗi $\rightarrow$ nhập lại $\rightarrow$ thành công & Luồng thay thế A1 & 047 \\ \hline
P3 & Gửi dữ liệu trùng $\rightarrow$ báo lỗi $\rightarrow$ đổi dữ liệu $\rightarrow$ thành công & Luồng thay thế A2 & 048 \\ \hline
P4 & Nhập dữ liệu $\rightarrow$ Cancel $\rightarrow$ không tạo user & Luồng thay thế A4 & 049 \\ \hline
P5 & Chọn Generate password $\rightarrow$ Create user $\rightarrow$ thành công & Luồng thay thế A3 & 050 \\ \hline
\end{tabular}
\caption{Scenario và test case của use-case testing}
\end{table}

\subsection{Danh sách test case của Feature 001}
Tổng cộng \textbf{50 test case}: ECP 001--017, BVA 018--040, Decision table 041--045, Use-case 046--050.
Các test case ghi ``ECP/BVA'' vừa là lớp không hợp lệ vừa là giá trị biên.
Kết quả mong đợi chung: với test case hợp lệ -- chuyển về Browse list of users, hiển thị ``Changes saved'' và user mới có trong danh sách;
với test case không hợp lệ -- form báo lỗi tại trường tương ứng (thông báo lấy từ mã nguồn Moodle) và user không được tạo.
Với các test case vượt giới hạn độ dài, thao tác kiểm tra là ô nhập chỉ nhận tối đa đúng số ký tự cho phép
(TC-001-008, 010, 012, 016, 017); riêng TC-001-003 (username 101 ký tự).

{\footnotesize
\begin{longtable}{|l|L{4.4cm}|c|L{6.4cm}|}
\hline
\rowcolor{gray!15}\textbf{ID} & \textbf{Tên test case} & \textbf{Kỹ thuật} & \textbf{Mô tả} \\ \hline
\endhead
TC-001-001 & Tạo user với dữ liệu hợp lệ & ECP & Mọi trường thuộc lớp hợp lệ V1-V6 (có nhập cả City). \\ \hline
TC-001-002 & Username rỗng & ECP & Lớp không hợp lệ I1: để trống Username. \\ \hline
TC-001-003 & Username 101 ký tự & ECP/BVA & Lớp không hợp lệ I2: Username dài 101 ký tự (cột username tối đa 100). \\ \hline
TC-001-004 & Username chứa chữ hoa & ECP & Lớp không hợp lệ I3: Username có chữ in hoa. \\ \hline
TC-001-005 & Username chứa ký tự không cho phép & ECP & Lớp không hợp lệ I4: Username có khoảng trắng và dấu \#. \\ \hline
TC-001-006 & Username đã tồn tại & ECP & Lớp không hợp lệ I5: Username trùng user admin có sẵn. \\ \hline
TC-001-007 & Password rỗng & ECP & Lớp không hợp lệ I6: không chọn Generate password và để trống New password. \\ \hline
TC-001-008 & Password 129 ký tự & ECP/BVA & Lớp không hợp lệ I7: Password dài hơn 128 ký tự. \\ \hline
TC-001-009 & First name rỗng & ECP/BVA & Lớp không hợp lệ I8. \\ \hline
TC-001-010 & First name 101 ký tự & ECP/BVA & Lớp không hợp lệ I9: First name dài hơn 100 ký tự. \\ \hline
TC-001-011 & Surname rỗng & ECP/BVA & Lớp không hợp lệ I10. \\ \hline
TC-001-012 & Surname 101 ký tự & ECP/BVA & Lớp không hợp lệ I11: Surname dài hơn 100 ký tự. \\ \hline
TC-001-013 & Email rỗng & ECP/BVA & Lớp không hợp lệ I12. \\ \hline
TC-001-014 & Email sai định dạng & ECP & Lớp không hợp lệ I13: email không có ký tự @. \\ \hline
TC-001-015 & Email đã tồn tại & ECP & Lớp không hợp lệ I14: email trùng email của user admin. \\ \hline
TC-001-016 & Email 101 ký tự & ECP/BVA & Lớp không hợp lệ I15: Email dài hơn 100 ký tự. \\ \hline
TC-001-017 & City 121 ký tự & ECP/BVA & Lớp không hợp lệ I16: City dài hơn 120 ký tự. \\ \hline
TC-001-018 & Username 1 ký tự & BVA & Biên hợp lệ của Username: độ dài 1. \\ \hline
TC-001-019 & Username 2 ký tự & BVA & Biên hợp lệ của Username: độ dài 2. \\ \hline
TC-001-020 & Username 99 ký tự & BVA & Biên hợp lệ của Username: độ dài 99. \\ \hline
TC-001-021 & Username 100 ký tự & BVA & Biên hợp lệ của Username: độ dài 100. \\ \hline
TC-001-022 & Password 1 ký tự & BVA & Biên hợp lệ của Password: độ dài 1. \\ \hline
TC-001-023 & Password 2 ký tự & BVA & Biên hợp lệ của Password: độ dài 2. \\ \hline
TC-001-024 & Password 127 ký tự & BVA & Biên hợp lệ của Password: độ dài 127. \\ \hline
TC-001-025 & Password 128 ký tự & BVA & Biên hợp lệ của Password: độ dài 128. \\ \hline
TC-001-026 & First name 1 ký tự & BVA & Biên hợp lệ của First name: độ dài 1. \\ \hline
TC-001-027 & First name 2 ký tự & BVA & Biên hợp lệ của First name: độ dài 2. \\ \hline
TC-001-028 & First name 99 ký tự & BVA & Biên hợp lệ của First name: độ dài 99. \\ \hline
TC-001-029 & First name 100 ký tự & BVA & Biên hợp lệ của First name: độ dài 100. \\ \hline
TC-001-030 & Surname 1 ký tự & BVA & Biên hợp lệ của Surname: độ dài 1. \\ \hline
TC-001-031 & Surname 2 ký tự & BVA & Biên hợp lệ của Surname: độ dài 2. \\ \hline
TC-001-032 & Surname 99 ký tự & BVA & Biên hợp lệ của Surname: độ dài 99. \\ \hline
TC-001-033 & Surname 100 ký tự & BVA & Biên hợp lệ của Surname: độ dài 100. \\ \hline
TC-001-034 & Email hợp lệ ngắn nhất & BVA & Biên min của Email: a@b.co (6 ký tự). \\ \hline
TC-001-035 & Email 99 ký tự & BVA & Biên max-1 của Email. \\ \hline
TC-001-036 & Email 100 ký tự & BVA & Biên max của Email. \\ \hline
TC-001-037 & City 0 ký tự & BVA & Biên hợp lệ của City: độ dài 0. \\ \hline
TC-001-038 & City 1 ký tự & BVA & Biên hợp lệ của City: độ dài 1. \\ \hline
TC-001-039 & City 119 ký tự & BVA & Biên hợp lệ của City: độ dài 119. \\ \hline
TC-001-040 & City 120 ký tự & BVA & Biên hợp lệ của City: độ dài 120. \\ \hline
TC-001-041 & R1: Generate password & DT & Chọn Generate password and notify user, không nhập password. \\ \hline
TC-001-042 & R2: Generate password + Suspended & DT & Chọn Generate password và Suspended cùng lúc. \\ \hline
TC-001-043 & R4: Force password change & DT & Nhập password, tick Force password change. \\ \hline
TC-001-044 & R5: Suspended account & DT & Nhập password, tick Suspended. \\ \hline
TC-001-045 & R8: Generate password + username trùng & DT & Chọn Generate password và dùng username admin đã tồn tại. \\ \hline
TC-001-046 & UC-P1: Luồng chính tạo user & UC & Từ trang Browse list of users bấm Add a new user, tạo user và kiểm tra trong danh sách. \\ \hline
TC-001-047 & UC-P2: Bỏ trống rồi nhập lại & UC & Gửi form rỗng, nhận lỗi, nhập đủ dữ liệu rồi gửi lại. \\ \hline
TC-001-048 & UC-P3: Username trùng rồi đổi & UC & Gửi username trùng admin, nhận lỗi, đổi username rồi gửi lại. \\ \hline
TC-001-049 & UC-P4: Hủy tạo user & UC & Nhập dữ liệu rồi bấm Cancel. \\ \hline
TC-001-050 & UC-P5: Tạo user bằng Generate password & UC & Luồng thay thế: tick Generate password, ô New password bị khóa. \\ \hline
\end{longtable}
}


\subsection{Test case mẫu (đầy đủ bước và kết quả mong đợi)}
Sáu test case dưới đây minh họa cách ghi trong file \texttt{.xlsx}; các test case còn lại viết theo cùng khuôn.
Dữ liệu dùng tiền tố \texttt{tc001\_} để dễ dọn dẹp.

\bigskip\noindent\textbf{TC-001-001 -- Tạo user với dữ liệu hợp lệ}\\
\noindent\textit{Precondition:} Moodle demo (Moodle 5.1, Docker) đang chạy tại http:/\allowbreak{}/\allowbreak{}127.0.0.1:9501; ngôn ngữ giao diện Moodle và trình duyệt là English; đã đăng nhập bằng tài khoản admin (admin/\allowbreak{}admin, mật khẩu không bị thay đổi); chính sách mật khẩu của demo đang tắt (passwordpolicy = 0). Chưa tồn tại user có username/\allowbreak{}email dùng trong test.
\smallskip

{\footnotesize
\begin{longtable}{|c|L{3.0cm}|L{5.1cm}|L{5.6cm}|}
\hline
\rowcolor{gray!15}\textbf{\#} & \textbf{Bước} & \textbf{Mô tả} & \textbf{Kết quả mong đợi} \\ \hline
\endhead
1 & Mở trang Add a new user & Vào Site administration \textgreater{} Users \textgreater{} Accounts \textgreater{} Add a new user (hoặc mở \textless{}BASE\_\allowbreak{}URL\textgreater{}/\allowbreak{}user/\allowbreak{}editadvanced.php?id=-1). & Form Add a new user hiển thị với các ô Username, New password, First name, Surname, Email address. \\ \hline
2 & Nhập Username & Nhập tc001\_\allowbreak{}valid vào ô Username. & Ô Username hiển thị đúng giá trị đã nhập. \\ \hline
3 & Nhập New password & Nhập Test@1234 vào ô New password. & Ô New password hiển thị đúng giá trị đã nhập. \\ \hline
4 & Nhập First name & Nhập Nguyen vào ô First name. & Ô First name hiển thị đúng giá trị đã nhập. \\ \hline
5 & Nhập Surname & Nhập Van A vào ô Surname. & Ô Surname hiển thị đúng giá trị đã nhập. \\ \hline
6 & Nhập Email address & Nhập tc001\_\allowbreak{}valid@example.com vào ô Email address. & Ô Email address hiển thị đúng giá trị đã nhập. \\ \hline
7 & Nhập City/\allowbreak{}town & Nhập Ho Chi Minh vào ô City/\allowbreak{}town. & Ô City/\allowbreak{}town hiển thị đúng giá trị đã nhập. \\ \hline
8 & Bấm Create user & Bấm nút Create user. & -- \\ \hline
9 & Kiểm tra kết quả & Quan sát trang sau khi gửi. & Chuyển tới trang Browse list of users (/\allowbreak{}admin/\allowbreak{}user.php), hiển thị thông báo "Changes saved" và user mới (email tc001\_\allowbreak{}valid@example.com) có trong danh sách. \\ \hline
\end{longtable}
}

\bigskip\noindent\textbf{TC-001-004 -- Username chứa chữ hoa}\\
\noindent\textit{Precondition:} Moodle demo (Moodle 5.1, Docker) đang chạy tại http:/\allowbreak{}/\allowbreak{}127.0.0.1:9501; ngôn ngữ giao diện Moodle và trình duyệt là English; đã đăng nhập bằng tài khoản admin (admin/\allowbreak{}admin, mật khẩu không bị thay đổi); chính sách mật khẩu của demo đang tắt (passwordpolicy = 0). Chưa tồn tại user có username/\allowbreak{}email dùng trong test.
\smallskip

{\footnotesize
\begin{longtable}{|c|L{3.0cm}|L{5.1cm}|L{5.6cm}|}
\hline
\rowcolor{gray!15}\textbf{\#} & \textbf{Bước} & \textbf{Mô tả} & \textbf{Kết quả mong đợi} \\ \hline
\endhead
1 & Mở trang Add a new user & Vào Site administration \textgreater{} Users \textgreater{} Accounts \textgreater{} Add a new user (hoặc mở \textless{}BASE\_\allowbreak{}URL\textgreater{}/\allowbreak{}user/\allowbreak{}editadvanced.php?id=-1). & Form Add a new user hiển thị với các ô Username, New password, First name, Surname, Email address. \\ \hline
2 & Nhập Username & Nhập TC001Upper vào ô Username. & Ô Username hiển thị đúng giá trị đã nhập. \\ \hline
3 & Nhập New password & Nhập Test@1234 vào ô New password. & Ô New password hiển thị đúng giá trị đã nhập. \\ \hline
4 & Nhập First name & Nhập Nguyen vào ô First name. & Ô First name hiển thị đúng giá trị đã nhập. \\ \hline
5 & Nhập Surname & Nhập Van A vào ô Surname. & Ô Surname hiển thị đúng giá trị đã nhập. \\ \hline
6 & Nhập Email address & Nhập tc001\_\allowbreak{}un\_\allowbreak{}upper@example.com vào ô Email address. & Ô Email address hiển thị đúng giá trị đã nhập. \\ \hline
7 & Bấm Create user & Bấm nút Create user. & -- \\ \hline
8 & Kiểm tra lỗi & Quan sát form sau khi gửi. & Không tạo user; vẫn ở trang Add a new user và hiển thị lỗi "The username must be in lower case" dưới ô Username. \\ \hline
\end{longtable}
}

\bigskip\noindent\textbf{TC-001-008 -- Password 129 ký tự}\\
\noindent\textit{Precondition:} Moodle demo (Moodle 5.1, Docker) đang chạy tại http:/\allowbreak{}/\allowbreak{}127.0.0.1:9501; ngôn ngữ giao diện Moodle và trình duyệt là English; đã đăng nhập bằng tài khoản admin (admin/\allowbreak{}admin, mật khẩu không bị thay đổi); chính sách mật khẩu của demo đang tắt (passwordpolicy = 0). Chưa tồn tại user có username/\allowbreak{}email dùng trong test.
\smallskip

{\footnotesize
\begin{longtable}{|c|L{3.0cm}|L{5.1cm}|L{5.6cm}|}
\hline
\rowcolor{gray!15}\textbf{\#} & \textbf{Bước} & \textbf{Mô tả} & \textbf{Kết quả mong đợi} \\ \hline
\endhead
1 & Mở trang Add a new user & Vào Site administration \textgreater{} Users \textgreater{} Accounts \textgreater{} Add a new user (hoặc mở \textless{}BASE\_\allowbreak{}URL\textgreater{}/\allowbreak{}user/\allowbreak{}editadvanced.php?id=-1). & Form Add a new user hiển thị với các ô Username, New password, First name, Surname, Email address. \\ \hline
2 & Nhập chuỗi vượt giới hạn & Gõ hoặc dán chuỗi 129 ký tự gồm 129 chữ "a" vào ô New password. & Ô New password chỉ nhận tối đa 128 ký tự; ký tự thứ 129 không nhập được. \\ \hline
3 & Kiểm tra độ dài & Đếm số ký tự đang có trong ô New password. & Ô chứa đúng 128 ký tự (thuộc tính maxlength = 128). \\ \hline
\end{longtable}
}

\bigskip\noindent\textbf{TC-001-021 -- Username 100 ký tự}\\
\noindent\textit{Precondition:} Moodle demo (Moodle 5.1, Docker) đang chạy tại http:/\allowbreak{}/\allowbreak{}127.0.0.1:9501; ngôn ngữ giao diện Moodle và trình duyệt là English; đã đăng nhập bằng tài khoản admin (admin/\allowbreak{}admin, mật khẩu không bị thay đổi); chính sách mật khẩu của demo đang tắt (passwordpolicy = 0). Chưa tồn tại user có username/\allowbreak{}email dùng trong test.
\smallskip

{\footnotesize
\begin{longtable}{|c|L{3.0cm}|L{5.1cm}|L{5.6cm}|}
\hline
\rowcolor{gray!15}\textbf{\#} & \textbf{Bước} & \textbf{Mô tả} & \textbf{Kết quả mong đợi} \\ \hline
\endhead
1 & Mở trang Add a new user & Vào Site administration \textgreater{} Users \textgreater{} Accounts \textgreater{} Add a new user (hoặc mở \textless{}BASE\_\allowbreak{}URL\textgreater{}/\allowbreak{}user/\allowbreak{}editadvanced.php?id=-1). & Form Add a new user hiển thị với các ô Username, New password, First name, Surname, Email address. \\ \hline
2 & Nhập Username & Nhập chuỗi 100 ký tự gồm "u" + 99 chữ "a" vào ô Username. & Ô Username chứa giá trị vừa nhập. \\ \hline
3 & Nhập New password & Nhập Test@1234 vào ô New password. & Ô New password hiển thị đúng giá trị đã nhập. \\ \hline
4 & Nhập First name & Nhập Nguyen vào ô First name. & Ô First name hiển thị đúng giá trị đã nhập. \\ \hline
5 & Nhập Surname & Nhập Van A vào ô Surname. & Ô Surname hiển thị đúng giá trị đã nhập. \\ \hline
6 & Nhập Email address & Nhập tc001\_\allowbreak{}un100@example.com vào ô Email address. & Ô Email address hiển thị đúng giá trị đã nhập. \\ \hline
7 & Bấm Create user & Bấm nút Create user. & -- \\ \hline
8 & Kiểm tra kết quả & Quan sát trang sau khi gửi. & Chuyển tới trang Browse list of users (/\allowbreak{}admin/\allowbreak{}user.php), hiển thị thông báo "Changes saved" và user mới (email tc001\_\allowbreak{}un100@example.com) có trong danh sách. \\ \hline
\end{longtable}
}

\bigskip\noindent\textbf{TC-001-043 -- R4: Force password change}\\
\noindent\textit{Precondition:} Moodle demo (Moodle 5.1, Docker) đang chạy tại http:/\allowbreak{}/\allowbreak{}127.0.0.1:9501; ngôn ngữ giao diện Moodle và trình duyệt là English; đã đăng nhập bằng tài khoản admin (admin/\allowbreak{}admin, mật khẩu không bị thay đổi); chính sách mật khẩu của demo đang tắt (passwordpolicy = 0). Chưa tồn tại user có username/\allowbreak{}email dùng trong test.
\smallskip

{\footnotesize
\begin{longtable}{|c|L{3.0cm}|L{5.1cm}|L{5.6cm}|}
\hline
\rowcolor{gray!15}\textbf{\#} & \textbf{Bước} & \textbf{Mô tả} & \textbf{Kết quả mong đợi} \\ \hline
\endhead
1 & Mở trang Add a new user & Vào Site administration \textgreater{} Users \textgreater{} Accounts \textgreater{} Add a new user (hoặc mở \textless{}BASE\_\allowbreak{}URL\textgreater{}/\allowbreak{}user/\allowbreak{}editadvanced.php?id=-1). & Form Add a new user hiển thị với các ô Username, New password, First name, Surname, Email address. \\ \hline
2 & Nhập Username & Nhập tc001\_\allowbreak{}dt\_\allowbreak{}r4 vào ô Username. & Ô Username hiển thị đúng giá trị đã nhập. \\ \hline
3 & Nhập New password & Nhập Test@1234 vào ô New password. & Ô New password hiển thị đúng giá trị đã nhập. \\ \hline
4 & Nhập First name & Nhập Nguyen vào ô First name. & Ô First name hiển thị đúng giá trị đã nhập. \\ \hline
5 & Nhập Surname & Nhập ForceR vào ô Surname. & Ô Surname hiển thị đúng giá trị đã nhập. \\ \hline
6 & Nhập Email address & Nhập tc001\_\allowbreak{}dt\_\allowbreak{}r4@example.com vào ô Email address. & Ô Email address hiển thị đúng giá trị đã nhập. \\ \hline
7 & Chọn Force password change & Tick ô Force password change. & Ô Force password change được chọn. \\ \hline
8 & Bấm Create user & Bấm nút Create user. & -- \\ \hline
9 & Kiểm tra kết quả & Quan sát trang sau khi gửi. & Chuyển tới trang Browse list of users (/\allowbreak{}admin/\allowbreak{}user.php), hiển thị thông báo "Changes saved" và user mới (email tc001\_\allowbreak{}dt\_\allowbreak{}r4@example.com) có trong danh sách. \\ \hline
10 & Mở trang Edit profile của user & Bấm vào tên user vừa tạo, rồi bấm Edit profile. & Trang chỉnh sửa user hiển thị. \\ \hline
11 & Kiểm tra cờ & Quan sát ô Force password change. & Ô Force password change đang được chọn. \\ \hline
\end{longtable}
}

\bigskip\noindent\textbf{TC-001-047 -- UC-P2: Bỏ trống rồi nhập lại}\\
\noindent\textit{Precondition:} Moodle demo (Moodle 5.1, Docker) đang chạy tại http:/\allowbreak{}/\allowbreak{}127.0.0.1:9501; ngôn ngữ giao diện Moodle và trình duyệt là English; đã đăng nhập bằng tài khoản admin (admin/\allowbreak{}admin, mật khẩu không bị thay đổi); chính sách mật khẩu của demo đang tắt (passwordpolicy = 0). Chưa tồn tại user có username/\allowbreak{}email dùng trong test.
\smallskip

{\footnotesize
\begin{longtable}{|c|L{3.0cm}|L{5.1cm}|L{5.6cm}|}
\hline
\rowcolor{gray!15}\textbf{\#} & \textbf{Bước} & \textbf{Mô tả} & \textbf{Kết quả mong đợi} \\ \hline
\endhead
1 & Mở trang Add a new user & Vào Site administration \textgreater{} Users \textgreater{} Accounts \textgreater{} Add a new user (hoặc mở \textless{}BASE\_\allowbreak{}URL\textgreater{}/\allowbreak{}user/\allowbreak{}editadvanced.php?id=-1). & Form Add a new user hiển thị với các ô Username, New password, First name, Surname, Email address. \\ \hline
2 & Gửi form rỗng & Không nhập gì, bấm Create user. & -- \\ \hline
3 & Kiểm tra lỗi lần 1 & Quan sát form. & Form không được gửi; hiển thị lỗi "Missing given name", "Missing last name" và "Required" (Email address). \\ \hline
4 & Nhập Username & Nhập tc001\_\allowbreak{}uc2 vào ô Username. & Ô Username hiển thị đúng giá trị đã nhập. \\ \hline
5 & Nhập New password & Nhập Test@1234 vào ô New password. & Ô New password hiển thị đúng giá trị đã nhập. \\ \hline
6 & Nhập First name & Nhập Nguyen vào ô First name. & Ô First name hiển thị đúng giá trị đã nhập. \\ \hline
7 & Nhập Surname & Nhập Van A vào ô Surname. & Ô Surname hiển thị đúng giá trị đã nhập. \\ \hline
8 & Nhập Email address & Nhập tc001\_\allowbreak{}uc2@example.com vào ô Email address. & Ô Email address hiển thị đúng giá trị đã nhập. \\ \hline
9 & Bấm Create user lần 2 & Bấm nút Create user. & -- \\ \hline
10 & Kiểm tra kết quả & Quan sát trang sau khi gửi. & Chuyển tới trang Browse list of users (/\allowbreak{}admin/\allowbreak{}user.php), hiển thị thông báo "Changes saved" và user mới (email tc001\_\allowbreak{}uc2@example.com) có trong danh sách. \\ \hline
\end{longtable}
}



% ================= 6 =================
% ================= Feature 002 (Lê Sĩ Nhật Quang) =================
\section{Feature 002: Add a new course}

\subsection{Mô tả và hướng dẫn sử dụng}
Quản trị viên tạo khóa học mới tại \textbf{Site administration $>$ Courses $>$ Add a new course}
(URL: \texttt{<BASE\_URL>/course/edit.php?category=0}), hoặc từ \textbf{Manage courses and categories $>$ Create new course}.
Form gồm các trường Course full name, Course short name, Course category, Course visibility, Course start date,
Course end date (có checkbox \textit{Enable}) và Course ID number; ngoài ra còn các nhóm thiết lập mở rộng
(Description, Course format, Appearance,\dots) giữ giá trị mặc định.

\textbf{Các bước sử dụng:}
\begin{enumerate}[noitemsep]
  \item Đăng nhập bằng tài khoản admin.
  \item Mở trang Add a new course.
  \item Nhập Course full name, Course short name; chọn Course category bằng ô tìm kiếm (autocomplete).
  \item Chọn Course visibility (Show/Hide), đặt Course start date; bật/tắt và đặt Course end date; nhập Course ID number (tùy chọn).
  \item Bấm \textbf{Save and display} (mở trang khóa học), \textbf{Save and return} (về trang quản lý) hoặc \textbf{Cancel}.
  \item Nếu dữ liệu không hợp lệ, form hiển thị lỗi tại trường tương ứng và không tạo khóa học.
\end{enumerate}

\subsection{Use case và activity diagram}
\begin{table}[H]
\centering\small
\begin{tabular}{|L{3.6cm}|L{11.4cm}|}
\hline
Tên & UC-002: Tạo khóa học mới \\ \hline
Tác nhân & Quản trị viên (Admin) \\ \hline
Tiền điều kiện & Admin đã đăng nhập; tồn tại ít nhất một category (``Category 1''). \\ \hline
Hậu điều kiện (thành công) & Khóa học mới được lưu trong category đã chọn và hiển thị trong trang quản lý khóa học. \\ \hline
Luồng chính & Mở form $\rightarrow$ nhập dữ liệu hợp lệ $\rightarrow$ Save and display $\rightarrow$ hệ thống tạo khóa học và mở trang khóa học. \\ \hline
Luồng thay thế & A1: bấm Save and return $\rightarrow$ tạo khóa học, quay về trang quản lý. A2: bấm Cancel $\rightarrow$ không tạo khóa học. \\ \hline
Luồng ngoại lệ & E1: thiếu trường bắt buộc (full name, short name, category). E2: short name đã tồn tại. E3: ID number đã tồn tại.
E4: end date không lớn hơn start date. E5: vượt độ dài tối đa. Mọi ngoại lệ: báo lỗi, giữ form để sửa và gửi lại. \\ \hline
\end{tabular}
\caption{Đặc tả use case UC-002}
\end{table}

\begin{figure}[H]
\centering
\begin{tikzpicture}[>=Stealth,
  act/.style={rectangle, rounded corners=4pt, draw, minimum width=3.6cm, minimum height=0.8cm, align=center, font=\small, fill=blue!5},
  dec/.style={diamond, aspect=2.2, draw, align=center, font=\scriptsize, inner sep=1pt, fill=yellow!15},
  st/.style={circle, fill=black, minimum size=0.35cm},
  en/.style={circle, draw, double, minimum size=0.35cm}]
\node[st] (s) {};
\node[act, below=0.5cm of s] (a1) {Mở trang Add a new course};
\node[act, below=0.5cm of a1] (a2) {Nhập/sửa dữ liệu form};
\node[dec, below=0.6cm of a2] (d0) {Bấm\\Cancel?};
\node[act, below=0.7cm of d0] (a3) {Bấm Save and display\\hoặc Save and return};
\node[dec, below=0.6cm of a3] (d1) {Đủ trường\\bắt buộc, đúng\\độ dài?};
\node[dec, below=0.7cm of d1] (d2) {Short name/\\ID number\\không trùng?};
\node[dec, below=0.7cm of d2] (d3) {End date tắt\\hoặc $>$ start?};
\node[dec, below=0.7cm of d3] (d4) {Save and\\display?};
\node[act, below left=0.7cm and -0.6cm of d4] (a4) {Mở trang khóa học};
\node[act, below right=0.7cm and -0.6cm of d4] (a5) {Về trang quản lý};
\node[en, below=2.2cm of d4] (e) {};
\node[act, right=2.3cm of d0] (c1) {Về trang quản lý,\\không tạo course};
\node[act, right=2.3cm of d1] (er1) {Lỗi Missing/\\Maximum of\dots};
\node[act, right=2.3cm of d2] (er2) {Lỗi ``already used''};
\node[act, right=2.3cm of d3] (er3) {Lỗi end date};
\draw[->] (s)--(a1); \draw[->] (a1)--(a2); \draw[->] (a2)--(d0);
\draw[->] (d0)--node[right,font=\scriptsize]{Không}(a3);
\draw[->] (d0)--node[above,font=\scriptsize]{Có}(c1);
\draw[->] (a3)--(d1);
\draw[->] (d1)--node[right,font=\scriptsize]{Có}(d2);
\draw[->] (d1)--node[above,font=\scriptsize]{Không}(er1);
\draw[->] (d2)--node[right,font=\scriptsize]{Có}(d3);
\draw[->] (d2)--node[above,font=\scriptsize]{Không}(er2);
\draw[->] (d3)--node[right,font=\scriptsize]{Có}(d4);
\draw[->] (d3)--node[above,font=\scriptsize]{Không}(er3);
\draw[->] (d4)-|node[above,font=\scriptsize,pos=0.25]{Có}(a4);
\draw[->] (d4)-|node[above,font=\scriptsize,pos=0.25]{Không}(a5);
\draw[->] (a4)|-(e); \draw[->] (a5)|-(e);
\coordinate (R) at ($(er1.east)+(0.5,0)$);
\coordinate (RC) at ($(R)+(0.7,0)$);
\draw (er1.east)--(er1.east -| R); \draw (er2.east)--(er2.east -| R); \draw (er3.east)--(er3.east -| R);
\draw[->] (er3.east -| R)--(a2.east -| R)--(a2.east);
\node[font=\scriptsize, fill=white, inner sep=1pt] at ($(a2.east -| R)!0.5!(a2.east)$) {sửa lỗi};
\draw[->] (c1.east)--(c1.east -| RC)|-(e.east);
\end{tikzpicture}
\caption{Activity diagram của UC-002 (các nhánh lỗi quay lại bước nhập dữ liệu; Cancel kết thúc không tạo khóa học)}
\end{figure}

\subsection{Các biến đầu vào}
\begin{table}[H]
\centering\small
\begin{tabular}{|L{2.8cm}|L{1.8cm}|L{5.0cm}|L{4.4cm}|}
\hline
\rowcolor{gray!15}\textbf{Biến} & \textbf{Bắt buộc?} & \textbf{Ràng buộc} & \textbf{Thông báo lỗi} \\ \hline
Course full name & Có & 1--254 ký tự; chuỗi chỉ có khoảng trắng coi như rỗng & Missing full name / Maximum of 254 characters \\ \hline
Course short name & Có & 1--100 ký tự; không trùng khóa học khác & Missing short name / Maximum of 100 characters / Short name is already used for another course (\dots) \\ \hline
Course category & Có & Chọn một category tồn tại & You must supply a value here. \\ \hline
Course visibility & Có (mặc định Show) & Show / Hide & Không có giá trị sai qua giao diện \\ \hline
Course start date & Có & Ngày giờ chọn từ dropdown (năm 1900--2050) & Không có giá trị sai qua giao diện \\ \hline
Course end date & Không (checkbox Enable) & Nếu Enable: phải $>$ start date & The course end date must be after the start date. \\ \hline
Course ID number & Không & 0--100 ký tự; nếu có thì không trùng & ID number is already used for another course (\dots) \\ \hline
\end{tabular}
\caption{Các biến đầu vào của Feature 002}
\end{table}
\textit{Ghi chú:} giới hạn độ dài trong bảng là giới hạn \emph{thiết kế} (đặc tả ban đầu: thuộc tính \texttt{maxlength} và quy tắc
kiểm tra phía trình duyệt của form). Khi chạy thực tế bằng Katalon Recorder (lệnh \texttt{type} gán giá trị trực tiếp cho ô nhập),
Moodle demo \textbf{chấp nhận} full name chỉ có khoảng trắng, full name 255 ký tự, short name 101 ký tự và ID number 101 ký tự;
xem mục~\ref{sec:f002-result}.

\subsection{Kỹ thuật 1: Phân lớp tương đương (ECP)}
\textbf{Cách áp dụng:} chia miền giá trị mỗi biến thành lớp hợp lệ (V) và không hợp lệ (I). Hai test case ``tất cả hợp lệ''
(TC-002-001, 002) phủ hết các lớp V (mỗi lựa chọn Show/Hide, end date bật/tắt, ID có/không); mỗi lớp I có một test case riêng,
các trường khác giữ hợp lệ.

\begin{table}[H]
\centering\footnotesize
\begin{tabular}{|l|L{5.0cm}|L{7.4cm}|}
\hline
\rowcolor{gray!15}\textbf{Biến} & \textbf{Lớp hợp lệ (TC)} & \textbf{Lớp không hợp lệ (TC)} \\ \hline
Full name & V1: 1--254 ký tự (001, 002) & I1: rỗng (003); I2: chỉ khoảng trắng (004); I3: 255 ký tự (005) \\ \hline
Short name & V2: 1--100 ký tự, chưa tồn tại (001, 002) & I4: rỗng (006); I5: 101 ký tự (007); I6: trùng EXIST01 (008) \\ \hline
Category & V3: chọn Category 1 (001, 002) & I7: No selection (009) \\ \hline
Visibility & V4: Show (001); V5: Hide (002) & -- \\ \hline
End date & V6: tắt Enable (002); V7: bật, end $>$ start (001) & I8: end = start (010); I9: end $<$ start (011) \\ \hline
ID number & V8: rỗng (002); V9: 1--100 ký tự, chưa tồn tại (001) & I10: 101 ký tự (012); I11: trùng ID-EXIST (013) \\ \hline
\end{tabular}
\caption{Lớp tương đương của Feature 002 (ID rút gọn: 003 nghĩa là TC-002-003)}
\end{table}
\textit{Sau khi chạy:} các lớp I2, I3, I5, I10 được thiết kế là \emph{không hợp lệ} nhưng hệ thống thực tế chấp nhận
(khóa học vẫn được tạo). Trong file \texttt{.xlsx}, mô tả của TC-002-004, 005, 007, 012 được cập nhật thành ``lớp hợp lệ''
và kết quả mong đợi được sửa theo hành vi thực tế; sự khác biệt này được ghi nhận là phát hiện kiểm thử.

\subsection{Kỹ thuật 2: Phân tích giá trị biên (BVA)}
\textbf{Cách áp dụng:} chọn min$-1$, min, min$+1$, max$-1$, max, max$+1$ cho mỗi biến có miền liên tục. Với ngày, ``độ dài''
là khoảng cách end -- start với bước nhỏ nhất là 1 phút (đơn vị nhỏ nhất của dropdown). Giá trị không hợp lệ trùng ECP được tham chiếu lại.

\begin{table}[H]
\centering\footnotesize
\begin{tabular}{|l|c|c|c|c|c|c|}
\hline
\rowcolor{gray!15}\textbf{Biến} & \textbf{min$-1$} & \textbf{min} & \textbf{min$+1$} & \textbf{max$-1$} & \textbf{max} & \textbf{max$+1$} \\ \hline
Full name (1--254) & 0 (003) & 1 (014) & 2 (015) & 253 (016) & 254 (017) & 255 (005) \\ \hline
Short name (1--100) & 0 (006) & 1 (018) & 2 (019) & 99 (020) & 100 (021) & 101 (007) \\ \hline
ID number (0--100) & -- & 0 (002) & 1 (022) \;/ 2 (023) & 99 (024) & 100 (025) & 101 (012) \\ \hline
End $-$ start & $-1$ phút (026) & 0 (010)$^{*}$ & $+1$ phút (027) & \multicolumn{3}{c|}{không có cận trên (năm tối đa 2050)} \\ \hline
\end{tabular}
\caption{Giá trị biên của Feature 002 ($^{*}$end = start là giá trị \emph{không hợp lệ}: biên hợp lệ nhỏ nhất là +1 phút;
max$+1$ của full name, short name, ID number thực tế được chấp nhận, xem mục~\ref{sec:f002-result})}
\end{table}

\subsection{Kỹ thuật 3: Bảng quyết định (Decision table)}
\textbf{Điều kiện:} C1 -- full name hợp lệ; C2 -- short name hợp lệ và không trùng; C3 -- đã chọn category;
C4 -- bật end date; C5 -- end date $>$ start date.
\textbf{Hành động:} A1 -- tạo khóa học; A2 -- lỗi full name; A3 -- lỗi short name; A4 -- lỗi category; A5 -- lỗi end date.
Các lỗi bắt buộc (A2--A4) được kiểm tra phía trình duyệt nên hiển thị đồng thời; lỗi trùng và lỗi ngày được kiểm tra phía server.

\begin{table}[H]
\centering\footnotesize
\begin{tabular}{|l|c|c|c|c|c|c|c|}
\hline
\rowcolor{gray!15}\textbf{Luật} & R1 & R2 & R3 & R4 & R5 & R6 & R7 \\ \hline
C1 Full name hợp lệ & Y & Y & Y & N & Y & Y & Y \\ \hline
C2 Short name hợp lệ & Y & Y & Y & N & N (trùng) & N (rỗng) & Y \\ \hline
C3 Có category & Y & Y & Y & N & Y & Y & N \\ \hline
C4 Bật end date & N & Y & Y & -- & Y & -- & -- \\ \hline
C5 End $>$ start & -- & Y & N & -- & N & -- & -- \\ \hline
A1 Tạo khóa học & X & X & & & & & \\ \hline
A2 Lỗi full name & & & & X & & & \\ \hline
A3 Lỗi short name & & & & X & X & X & \\ \hline
A4 Lỗi category & & & & X & & & X \\ \hline
A5 Lỗi end date & & & X & & X & & \\ \hline
\textbf{Test case} & 028 & 029 & 030 & 031 & 032 & 006 & 009 \\ \hline
\end{tabular}
\caption{Decision table của Feature 002 (R6, R7 tái sử dụng test case ECP)}
\end{table}

\subsection{Kỹ thuật 4: Use-case testing}
\begin{table}[H]
\centering\footnotesize
\begin{tabular}{|c|L{7.4cm}|L{3.2cm}|c|}
\hline
\rowcolor{gray!15}\textbf{Path} & \textbf{Đường đi trong activity diagram} & \textbf{Loại} & \textbf{Test case} \\ \hline
P1 & Mở form $\rightarrow$ nhập hợp lệ $\rightarrow$ Save and display $\rightarrow$ trang khóa học & Luồng chính & 033 \\ \hline
P2 & Mở form $\rightarrow$ nhập hợp lệ $\rightarrow$ Save and return $\rightarrow$ trang quản lý & Luồng thay thế A1 & 034 \\ \hline
P3 & Nhập dữ liệu $\rightarrow$ Cancel $\rightarrow$ không tạo khóa học & Luồng thay thế A2 & 035 \\ \hline
P4 & Gửi form rỗng $\rightarrow$ lỗi bắt buộc $\rightarrow$ nhập lại $\rightarrow$ thành công & Ngoại lệ E1 & 036 \\ \hline
P5 & Short name trùng $\rightarrow$ lỗi ``already used'' $\rightarrow$ đổi $\rightarrow$ thành công & Ngoại lệ E2 & 037 \\ \hline
P6 & End date $<$ start $\rightarrow$ lỗi end date $\rightarrow$ sửa $\rightarrow$ thành công & Ngoại lệ E4 & 038 \\ \hline
\end{tabular}
\caption{Scenario và test case use-case testing của Feature 002 (E3 đã phủ bởi TC-002-013; E5 được kiểm tra bởi TC-002-005, 007, 012 nhưng thực tế hệ thống không báo lỗi, xem mục~\ref{sec:f002-result})}
\end{table}

\subsection{Kết quả thực thi và phát hiện}\label{sec:f002-result}
Toàn bộ 38 test case đã được ghi và chạy bằng Katalon Recorder trên Moodle demo (file \texttt{F002-Add a new course.krecorder}).
Kết quả mong đợi trong file \texttt{.xlsx} được cập nhật theo hành vi thực tế quan sát được. Các test case có kết quả khác đặc tả ban đầu:

\begin{table}[H]
\centering\footnotesize
\begin{tabular}{|l|L{3.3cm}|L{4.3cm}|L{4.6cm}|}
\hline
\rowcolor{gray!15}\textbf{Test case} & \textbf{Đầu vào} & \textbf{Kết quả dự kiến ban đầu} & \textbf{Kết quả thực tế} \\ \hline
TC-002-004 & Full name = 3 khoảng trắng & Lỗi ``Missing full name'' & Khóa học được tạo với full name là khoảng trắng \\ \hline
TC-002-005 & Full name 255 ký tự & Lỗi ``Maximum of 254 characters'' hoặc bị cắt còn 254 & Khóa học được tạo với full name \texttt{Faaa\dots} \\ \hline
TC-002-007 & Short name 101 ký tự & Lỗi ``Maximum of 100 characters'' & Khóa học được tạo \\ \hline
TC-002-012 & ID number 101 ký tự & Lỗi ``Maximum of 100 characters'' hoặc bị cắt còn 100 & Khóa học được tạo \\ \hline
\end{tabular}
\caption{Khác biệt giữa đặc tả ban đầu và hành vi thực tế của Feature 002}
\end{table}

\textbf{Nhận xét:}
\begin{itemize}[noitemsep]
  \item Giới hạn độ dài của form chỉ được áp dụng ở phía trình duyệt (\texttt{maxlength}, kiểm tra JavaScript). Lệnh \texttt{type}
        của Katalon gán giá trị trực tiếp nên vượt qua được, và phía server không chặn vì cột dữ liệu trong CSDL dài hơn
        giới hạn hiển thị trên form.
  \item Full name chỉ có khoảng trắng vẫn được lưu: khóa học có tên ``trống'' gây khó nhận diện trong danh sách. Đây là điểm
        nên báo cáo như một lỗi tiềm ẩn (mức độ thấp).
  \item Các test case còn lại cho kết quả đúng như mong đợi: các lỗi bắt buộc (Missing full name, Missing short name, You must supply a value here.),
        lỗi trùng short name/ID number và lỗi end date đều xuất hiện đúng trường.
\end{itemize}

\subsection{Danh sách test case của Feature 002}
Tổng cộng \textbf{38 test case}: ECP 001--013, BVA 014--027, Decision table 028--032, Use-case 033--038.
Kết quả mong đợi chung: test case hợp lệ -- khóa học được tạo và hiển thị trong \textit{My courses} với đúng full name;
test case không hợp lệ -- vẫn ở form Add a new course, hiển thị lỗi tại trường tương ứng và không có khóa học mới.
Riêng TC-002-004, 005, 007, 012 có kết quả mong đợi theo hành vi thực tế (mục~\ref{sec:f002-result}).
Dữ liệu dùng tiền tố \texttt{TC002} để dễ dọn dẹp.

{\footnotesize
\begin{longtable}{|l|L{4.4cm}|c|L{6.4cm}|}
\hline
\rowcolor{gray!15}\textbf{ID} & \textbf{Tên test case} & \textbf{Kỹ thuật} & \textbf{Mô tả} \\ \hline
\endhead
TC-002-001 & Tạo course với dữ liệu hợp lệ (Show, có end date, có ID) & ECP & Tất cả các trường thuộc lớp hợp lệ V1, V2, V3, V4, V7, V9. \\ \hline
TC-002-002 & Tạo course Hide, tắt end date, không ID & ECP & Lớp hợp lệ V1, V2, V3, V5, V6, V8. \\ \hline
TC-002-003 & Full name rỗng & ECP & Lớp không hợp lệ I1: để trống Course full name. \\ \hline
TC-002-004 & Full name chỉ có khoảng trắng & ECP & Lớp hợp lệ I2: Course full name gồm 3 khoảng trắng. \\ \hline
TC-002-005 & Full name 255 ký tự & ECP/BVA & Lớp hợp lệ I3: Course full name dài 255 ký tự (tối đa 254). \\ \hline
TC-002-006 & Short name rỗng & ECP & Lớp không hợp lệ I4: để trống Course short name. \\ \hline
TC-002-007 & Short name 101 ký tự & ECP/BVA & Lớp hợp lệ I5: Course short name dài 101 ký tự. \\ \hline
TC-002-008 & Short name đã tồn tại & ECP & Lớp không hợp lệ I6: Course short name trùng "EXIST01" của course có sẵn. \\ \hline
TC-002-009 & Không chọn Course category & ECP & Lớp không hợp lệ I7: Course category để "No selection". \\ \hline
TC-002-010 & End date bằng start date & ECP/BVA & Lớp không hợp lệ I8: end date = start date (1 January 2026 08:00). \\ \hline
TC-002-011 & End date trước start date & ECP & Lớp không hợp lệ I9: end date (1 January 2025 08:00) \textless{} start date. \\ \hline
TC-002-012 & ID number 101 ký tự & ECP/BVA & Lớp hợp lệ I10: Course ID number dài 101 ký tự (tối đa 100). \\ \hline
TC-002-013 & ID number đã tồn tại & ECP & Lớp không hợp lệ I11: Course ID number trùng "ID-EXIST". \\ \hline
TC-002-014 & Full name 1 ký tự & BVA & Biên hợp lệ của Course full name: độ dài 1. \\ \hline
TC-002-015 & Full name 2 ký tự & BVA & Biên hợp lệ của Course full name: độ dài 2. \\ \hline
TC-002-016 & Full name 253 ký tự & BVA & Biên hợp lệ của Course full name: độ dài 253. \\ \hline
TC-002-017 & Full name 254 ký tự & BVA & Biên hợp lệ của Course full name: độ dài 254. \\ \hline
TC-002-018 & Short name 1 ký tự & BVA & Biên hợp lệ của Course short name: độ dài 1. \\ \hline
TC-002-019 & Short name 2 ký tự & BVA & Biên hợp lệ của Course short name: độ dài 2. \\ \hline
TC-002-020 & Short name 99 ký tự & BVA & Biên hợp lệ của Course short name: độ dài 99. \\ \hline
TC-002-021 & Short name 100 ký tự & BVA & Biên hợp lệ của Course short name: độ dài 100. \\ \hline
TC-002-022 & ID number 1 ký tự & BVA & Biên hợp lệ của Course ID number: độ dài 1. \\ \hline
TC-002-023 & ID number 2 ký tự & BVA & Biên hợp lệ của Course ID number: độ dài 2. \\ \hline
TC-002-024 & ID number 99 ký tự & BVA & Biên hợp lệ của Course ID number: độ dài 99. \\ \hline
TC-002-025 & ID number 100 ký tự & BVA & Biên hợp lệ của Course ID number: độ dài 100. \\ \hline
TC-002-026 & End date = start date - 1 phút & BVA & Biên dưới không hợp lệ: end 1 January 2026 07:59, start 1 January 2026 08:00. \\ \hline
TC-002-027 & End date = start date + 1 phút & BVA & Biên hợp lệ nhỏ nhất: end 1 January 2026 08:01, start 1 January 2026 08:00. \\ \hline
TC-002-028 & DT-R1: Hợp lệ, end date tắt & DT & R1: C1=T, C2=T, C3=T, C4=F (C5 don't care) -\textgreater{} tạo course thành công. \\ \hline
TC-002-029 & DT-R2: Hợp lệ, end date \textgreater{} start & DT & R2: C1=T, C2=T, C3=T, C4=T, C5=T -\textgreater{} tạo course thành công. \\ \hline
TC-002-030 & DT-R3: End date không hợp lệ & DT & R3: C1=T, C2=T, C3=T, C4=T, C5=F -\textgreater{} lỗi end date. \\ \hline
TC-002-031 & DT-R4: Thiếu toàn bộ trường bắt buộc & DT & R4: C1=F, C2=F, C3=F -\textgreater{} hiển thị đồng thời 3 lỗi. \\ \hline
TC-002-032 & DT-R5: Short name trùng và end date sai & DT & R5: C1=T, C2=F (trùng), C3=T, C4=T, C5=F -\textgreater{} hiển thị cả 2 lỗi server. \\ \hline
TC-002-033 & UC-P1: Luồng chính - Save and display & UC & Basic flow: nhập dữ liệu hợp lệ, Save and display. \\ \hline
TC-002-034 & UC-P2: Luồng thay thế - Save and return & UC & Alternative flow A1: lưu và quay về trang quản lý. \\ \hline
TC-002-035 & UC-P3: Hủy tạo course & UC & Alternative flow A2: nhập dữ liệu rồi Cancel. \\ \hline
TC-002-036 & UC-P4: Bỏ trống rồi nhập lại & UC & Exception E1 -\textgreater{} sửa -\textgreater{} thành công. \\ \hline
TC-002-037 & UC-P5: Short name trùng rồi đổi & UC & Exception E2 -\textgreater{} sửa -\textgreater{} thành công. \\ \hline
TC-002-038 & UC-P6: End date sai rồi sửa & UC & Exception E4 -\textgreater{} sửa -\textgreater{} thành công. \\ \hline
\end{longtable}
}


\subsection{Test case mẫu của Feature 002}
Sáu test case dưới đây được trích nguyên từ file \texttt{.xlsx} (sheet \textit{Add a new course}); các test case còn lại cùng khuôn.

\bigskip\noindent\textbf{TC-002-001 -- Tạo course với dữ liệu hợp lệ (Show, có end date, có ID)}\\
\noindent\textit{Precondition:} Moodle demo (Moodle 5.1, Docker) đang chạy tại http:/\allowbreak{}/\allowbreak{}127.0.0.1:9501 giao diện Moodle và trình duyệt là English; đã đăng nhập bằng admin (admin/\allowbreak{}admin). Đã có category "Category 1". Script Setup đã tạo sẵn course Full name "Existing Course", Short name "EXIST01", ID number "ID-EXIST" trong "Category 1". Chưa tồn tại course nào có Short name/\allowbreak{}ID number dùng trong test case này (script Cleanup xóa các course TC002-* sau khi chạy).
\smallskip

{\footnotesize
\begin{longtable}{|c|L{3.0cm}|L{5.1cm}|L{5.6cm}|}
\hline
\rowcolor{gray!15}\textbf{\#} & \textbf{Bước} & \textbf{Mô tả} & \textbf{Kết quả mong đợi} \\ \hline
\endhead
1 & Mở trang Add a new course & Vào Site administration \textgreater{} Courses \textgreater{} Add a new course (hoặc mở \textless{}BASE\_\allowbreak{}URL\textgreater{}/\allowbreak{}course/\allowbreak{}edit.php?category=0). & Form "Add a new course" hiển thị: Course full name, Course short name, Course category (No selection), Course visibility, Course start date, Course end date, Course ID number. \\ \hline
2 & Nhập Course full name & Nhập "TC002 Course 001" vào ô Course full name. & Ô Course full name hiển thị giá trị đã nhập. \\ \hline
3 & Nhập Course short name & Nhập "TC002-001" vào ô Course short name. & Ô Course short name hiển thị giá trị đã nhập. \\ \hline
4 & Chọn Course category & Click ô Search của Course category, gõ "Category 1" và chọn "Category 1" trong danh sách. & Course category hiển thị "Category 1". \\ \hline
5 & Chọn Course visibility & Chọn "Show" trong dropdown Course visibility. & Dropdown hiển thị "Show". \\ \hline
6 & Đặt Course start date & Chọn ngày/\allowbreak{}tháng/\allowbreak{}năm/\allowbreak{}giờ/\allowbreak{}phút = 1 January 2026 08:00. & Course start date hiển thị 1 January 2026 08:00. \\ \hline
7 & Đặt Course end date & Tick "Enable" của Course end date, chọn 1 January 2027 08:00. & Course end date hiển thị 1 January 2027 08:00. \\ \hline
8 & Nhập Course ID number & Nhập "ID-002-001" vào ô Course ID number. & Ô Course ID number hiển thị giá trị đã nhập. \\ \hline
9 & Lưu course & Click nút "Save and display". & Không có thông báo lỗi; course được tạo và trang course hiển thị tiêu đề "TC002 Course 001"; breadcrumb chứa short name "TC002-001". \\ \hline
\end{longtable}
}

\bigskip\noindent\textbf{TC-002-004 -- Full name chỉ có khoảng trắng}\\
\noindent\textit{Precondition:} Moodle demo (Moodle 5.1, Docker) đang chạy tại http:/\allowbreak{}/\allowbreak{}127.0.0.1:9501 giao diện Moodle và trình duyệt là English; đã đăng nhập bằng admin (admin/\allowbreak{}admin). Đã có category "Category 1". Script Setup đã tạo sẵn course Full name "Existing Course", Short name "EXIST01", ID number "ID-EXIST" trong "Category 1". Chưa tồn tại course nào có Short name/\allowbreak{}ID number dùng trong test case này (script Cleanup xóa các course TC002-* sau khi chạy).
\smallskip

{\footnotesize
\begin{longtable}{|c|L{3.0cm}|L{5.1cm}|L{5.6cm}|}
\hline
\rowcolor{gray!15}\textbf{\#} & \textbf{Bước} & \textbf{Mô tả} & \textbf{Kết quả mong đợi} \\ \hline
\endhead
1 & Mở trang Add a new course & Vào Site administration \textgreater{} Courses \textgreater{} Add a new course (hoặc mở \textless{}BASE\_\allowbreak{}URL\textgreater{}/\allowbreak{}course/\allowbreak{}edit.php?category=0). & Form "Add a new course" hiển thị: Course full name, Course short name, Course category (No selection), Course visibility, Course start date, Course end date, Course ID number. \\ \hline
2 & Nhập Course full name & Nhập 3 ký tự khoảng trắng vào ô Course full name. & Ô Course full name hiển thị giá trị đã nhập. \\ \hline
3 & Nhập Course short name & Nhập "TC002-004" vào ô Course short name. & Ô Course short name hiển thị giá trị đã nhập. \\ \hline
4 & Chọn Course category & Click ô Search của Course category, gõ "Category 1" và chọn "Category 1" trong danh sách. & Course category hiển thị "Category 1". \\ \hline
5 & Chọn Course visibility & Chọn "Show" trong dropdown Course visibility. & Dropdown hiển thị "Show". \\ \hline
6 & Đặt Course start date & Chọn ngày/\allowbreak{}tháng/\allowbreak{}năm/\allowbreak{}giờ/\allowbreak{}phút = 1 January 2026 08:00. & Course start date hiển thị 1 January 2026 08:00. \\ \hline
7 & Đặt Course end date & Tick "Enable" của Course end date, chọn 1 January 2027 08:00. & Course end date hiển thị 1 January 2027 08:00. \\ \hline
8 & Nhập Course ID number & Để trống (xóa hết nội dung) vào ô Course ID number. & Ô Course ID number hiển thị giá trị đã nhập. \\ \hline
9 & Lưu course & Click nút "Save and display". & Form được lưu với tên là 3 khoảng trắng \\ \hline
\end{longtable}
}

\bigskip\noindent\textbf{TC-002-008 -- Short name đã tồn tại}\\
\noindent\textit{Precondition:} Moodle demo (Moodle 5.1, Docker) đang chạy tại http:/\allowbreak{}/\allowbreak{}127.0.0.1:9501 giao diện Moodle và trình duyệt là English; đã đăng nhập bằng admin (admin/\allowbreak{}admin). Đã có category "Category 1". Script Setup đã tạo sẵn course Full name "Existing Course", Short name "EXIST01", ID number "ID-EXIST" trong "Category 1". Chưa tồn tại course nào có Short name/\allowbreak{}ID number dùng trong test case này (script Cleanup xóa các course TC002-* sau khi chạy).
\smallskip

{\footnotesize
\begin{longtable}{|c|L{3.0cm}|L{5.1cm}|L{5.6cm}|}
\hline
\rowcolor{gray!15}\textbf{\#} & \textbf{Bước} & \textbf{Mô tả} & \textbf{Kết quả mong đợi} \\ \hline
\endhead
1 & Mở trang Add a new course & Vào Site administration \textgreater{} Courses \textgreater{} Add a new course (hoặc mở \textless{}BASE\_\allowbreak{}URL\textgreater{}/\allowbreak{}course/\allowbreak{}edit.php?category=0). & Form "Add a new course" hiển thị: Course full name, Course short name, Course category (No selection), Course visibility, Course start date, Course end date, Course ID number. \\ \hline
2 & Nhập Course full name & Nhập "TC002 Course 008" vào ô Course full name. & Ô Course full name hiển thị giá trị đã nhập. \\ \hline
3 & Nhập Course short name & Nhập "EXIST01" vào ô Course short name. & Ô Course short name hiển thị giá trị đã nhập. \\ \hline
4 & Chọn Course category & Click ô Search của Course category, gõ "Category 1" và chọn "Category 1" trong danh sách. & Course category hiển thị "Category 1". \\ \hline
5 & Chọn Course visibility & Chọn "Show" trong dropdown Course visibility. & Dropdown hiển thị "Show". \\ \hline
6 & Đặt Course start date & Chọn ngày/\allowbreak{}tháng/\allowbreak{}năm/\allowbreak{}giờ/\allowbreak{}phút = 1 January 2026 08:00. & Course start date hiển thị 1 January 2026 08:00. \\ \hline
7 & Đặt Course end date & Tick "Enable" của Course end date, chọn 1 January 2027 08:00. & Course end date hiển thị 1 January 2027 08:00. \\ \hline
8 & Nhập Course ID number & Để trống (xóa hết nội dung) vào ô Course ID number. & Ô Course ID number hiển thị giá trị đã nhập. \\ \hline
9 & Lưu course & Click nút "Save and display". & Form không được lưu, vẫn ở trang Add a new course; hiển thị lỗi "Short name is already used for another course (Existing Course)". Không có course mới được tạo. \\ \hline
\end{longtable}
}

\bigskip\noindent\textbf{TC-002-010 -- End date bằng start date}\\
\noindent\textit{Precondition:} Moodle demo (Moodle 5.1, Docker) đang chạy tại http:/\allowbreak{}/\allowbreak{}127.0.0.1:9501 giao diện Moodle và trình duyệt là English; đã đăng nhập bằng admin (admin/\allowbreak{}admin). Đã có category "Category 1". Script Setup đã tạo sẵn course Full name "Existing Course", Short name "EXIST01", ID number "ID-EXIST" trong "Category 1". Chưa tồn tại course nào có Short name/\allowbreak{}ID number dùng trong test case này (script Cleanup xóa các course TC002-* sau khi chạy).
\smallskip

{\footnotesize
\begin{longtable}{|c|L{3.0cm}|L{5.1cm}|L{5.6cm}|}
\hline
\rowcolor{gray!15}\textbf{\#} & \textbf{Bước} & \textbf{Mô tả} & \textbf{Kết quả mong đợi} \\ \hline
\endhead
1 & Mở trang Add a new course & Vào Site administration \textgreater{} Courses \textgreater{} Add a new course (hoặc mở \textless{}BASE\_\allowbreak{}URL\textgreater{}/\allowbreak{}course/\allowbreak{}edit.php?category=0). & Form "Add a new course" hiển thị: Course full name, Course short name, Course category (No selection), Course visibility, Course start date, Course end date, Course ID number. \\ \hline
2 & Nhập Course full name & Nhập "TC002 Course 010" vào ô Course full name. & Ô Course full name hiển thị giá trị đã nhập. \\ \hline
3 & Nhập Course short name & Nhập "TC002-010" vào ô Course short name. & Ô Course short name hiển thị giá trị đã nhập. \\ \hline
4 & Chọn Course category & Click ô Search của Course category, gõ "Category 1" và chọn "Category 1" trong danh sách. & Course category hiển thị "Category 1". \\ \hline
5 & Chọn Course visibility & Chọn "Show" trong dropdown Course visibility. & Dropdown hiển thị "Show". \\ \hline
6 & Đặt Course start date & Chọn ngày/\allowbreak{}tháng/\allowbreak{}năm/\allowbreak{}giờ/\allowbreak{}phút = 1 January 2026 08:00. & Course start date hiển thị 1 January 2026 08:00. \\ \hline
7 & Đặt Course end date & Tick "Enable" của Course end date, chọn 1 January 2026 08:00. & Course end date hiển thị 1 January 2026 08:00. \\ \hline
8 & Nhập Course ID number & Để trống (xóa hết nội dung) vào ô Course ID number. & Ô Course ID number hiển thị giá trị đã nhập. \\ \hline
9 & Lưu course & Click nút "Save and display". & Form không được lưu, vẫn ở trang Add a new course; hiển thị lỗi "The course end date must be after the start date.". Không có course mới được tạo. \\ \hline
\end{longtable}
}

\bigskip\noindent\textbf{TC-002-031 -- DT-R4: Thiếu toàn bộ trường bắt buộc}\\
\noindent\textit{Precondition:} Moodle demo (Moodle 5.1, Docker) đang chạy tại http:/\allowbreak{}/\allowbreak{}127.0.0.1:9501 giao diện Moodle và trình duyệt là English; đã đăng nhập bằng admin (admin/\allowbreak{}admin). Đã có category "Category 1". Script Setup đã tạo sẵn course Full name "Existing Course", Short name "EXIST01", ID number "ID-EXIST" trong "Category 1". Chưa tồn tại course nào có Short name/\allowbreak{}ID number dùng trong test case này (script Cleanup xóa các course TC002-* sau khi chạy).
\smallskip

{\footnotesize
\begin{longtable}{|c|L{3.0cm}|L{5.1cm}|L{5.6cm}|}
\hline
\rowcolor{gray!15}\textbf{\#} & \textbf{Bước} & \textbf{Mô tả} & \textbf{Kết quả mong đợi} \\ \hline
\endhead
1 & Mở trang Add a new course & Vào Site administration \textgreater{} Courses \textgreater{} Add a new course (hoặc mở \textless{}BASE\_\allowbreak{}URL\textgreater{}/\allowbreak{}course/\allowbreak{}edit.php?category=0). & Form "Add a new course" hiển thị: Course full name, Course short name, Course category (No selection), Course visibility, Course start date, Course end date, Course ID number. \\ \hline
2 & Nhập Course full name & Để trống (xóa hết nội dung) vào ô Course full name. & Ô Course full name hiển thị giá trị đã nhập. \\ \hline
3 & Nhập Course short name & Để trống (xóa hết nội dung) vào ô Course short name. & Ô Course short name hiển thị giá trị đã nhập. \\ \hline
4 & Không chọn Course category & Giữ nguyên Course category là "No selection". & Course category vẫn là "No selection". \\ \hline
5 & Chọn Course visibility & Chọn "Show" trong dropdown Course visibility. & Dropdown hiển thị "Show". \\ \hline
6 & Đặt Course start date & Chọn ngày/\allowbreak{}tháng/\allowbreak{}năm/\allowbreak{}giờ/\allowbreak{}phút = 1 January 2026 08:00. & Course start date hiển thị 1 January 2026 08:00. \\ \hline
7 & Đặt Course end date & Tick "Enable" của Course end date, chọn 1 January 2027 08:00. & Course end date hiển thị 1 January 2027 08:00. \\ \hline
8 & Nhập Course ID number & Để trống (xóa hết nội dung) vào ô Course ID number. & Ô Course ID number hiển thị giá trị đã nhập. \\ \hline
9 & Lưu course & Click nút "Save and display". & Form không được lưu, vẫn ở trang Add a new course; hiển thị lỗi "Missing full name", "Missing short name" và "You must supply a value here." cùng lúc. Không có course mới được tạo. \\ \hline
\end{longtable}
}

\bigskip\noindent\textbf{TC-002-037 -- UC-P5: Short name trùng rồi đổi}\\
\noindent\textit{Precondition:} Moodle demo (Moodle 5.1, Docker) đang chạy tại http:/\allowbreak{}/\allowbreak{}127.0.0.1:9501 giao diện Moodle và trình duyệt là English; đã đăng nhập bằng admin (admin/\allowbreak{}admin). Đã có category "Category 1". Script Setup đã tạo sẵn course Full name "Existing Course", Short name "EXIST01", ID number "ID-EXIST" trong "Category 1". Chưa tồn tại course nào có Short name/\allowbreak{}ID number dùng trong test case này (script Cleanup xóa các course TC002-* sau khi chạy).
\smallskip

{\footnotesize
\begin{longtable}{|c|L{3.0cm}|L{5.1cm}|L{5.6cm}|}
\hline
\rowcolor{gray!15}\textbf{\#} & \textbf{Bước} & \textbf{Mô tả} & \textbf{Kết quả mong đợi} \\ \hline
\endhead
1 & Mở trang quản lý course & Vào Site administration \textgreater{} Courses \textgreater{} Manage courses and categories, chọn "Category 1". & Trang Course and category management hiển thị, liệt kê course của Category 1. \\ \hline
2 & Mở form tạo course & Click link "Create new course". & Form "Add a new course" hiển thị, Course category mặc định là "Category 1". \\ \hline
3 & Nhập Course full name & Nhập "TC002 UC Dup" vào ô Course full name. & Ô hiển thị giá trị đã nhập. \\ \hline
4 & Nhập Course short name & Nhập "EXIST01" vào ô Course short name. & Ô hiển thị giá trị đã nhập. \\ \hline
5 & Gửi form & Click "Save and display". & Hiển thị "Short name is already used for another course (Existing Course)". \\ \hline
6 & Đổi Short name & Xóa ô Course short name, nhập "TC002-UC5". & Ô hiển thị "TC002-UC5". \\ \hline
7 & Gửi lại & Click "Save and display". & Trang course "TC002 UC Dup" hiển thị. \\ \hline
\end{longtable}
}



\section{Feature 003: ---}
\section{Feature 004: Add new category}

\subsection{Mô tả và hướng dẫn sử dụng}
Quản trị viên tạo danh mục khóa học (\textit{course category}) mới tại \textbf{Site administration $>$ Courses $>$ Manage courses and categories $>$ Create new category}
(URL: \texttt{<BASE\_URL>/course/editcategory.php?parent=0}). Danh mục có thể nằm ở cấp Top hoặc là danh mục con của một danh mục có sẵn.

\textbf{Các bước sử dụng:}
\begin{enumerate}[noitemsep]
  \item Đăng nhập bằng tài khoản admin.
  \item Mở trang Manage courses and categories, bấm \textbf{Create new category}.
  \item Chọn \textit{Parent category} (Top hoặc một danh mục), nhập \textit{Category name}; có thể nhập thêm \textit{Category ID number} và \textit{Description}.
  \item Bấm \textbf{Create category} (hoặc \textbf{Cancel} để hủy).
  \item Hệ thống kiểm tra dữ liệu; nếu hợp lệ thì tạo danh mục và chuyển về trang Manage courses and categories, nếu không thì hiển thị lỗi và
        giữ nguyên form (riêng Name quá dài sẽ hiện trang lỗi riêng của Moodle).
\end{enumerate}

\subsection{Use case và activity diagram}
\textbf{Use case:}
\begin{table}[H]
\centering\small
\begin{tabular}{|L{3.6cm}|L{11.4cm}|}
\hline
Tên & UC-004: Tạo danh mục khóa học mới \\ \hline
Tác nhân & Quản trị viên (Admin) \\ \hline
Tiền điều kiện & Admin đã đăng nhập; có quyền quản lý danh mục khóa học. \\ \hline
Hậu điều kiện (thành công) & Danh mục mới được lưu và hiển thị trong Manage courses and categories, đúng vị trí dưới Parent đã chọn. \\ \hline
Luồng chính & Mở form $\rightarrow$ chọn Parent, nhập Name $\rightarrow$ Create category $\rightarrow$ hệ thống tạo danh mục và chuyển về trang quản lý. \\ \hline
Luồng thay thế & A1: thiếu Name hoặc chưa chọn Parent $\rightarrow$ báo lỗi, nhập lại. A2: ID number trùng $\rightarrow$ báo lỗi, đổi ID.
A3: Name quá 255 ký tự $\rightarrow$ trang lỗi, quay lại form. A4: chọn danh mục làm Parent để tạo danh mục con. A5: bấm Cancel, không tạo danh mục. \\ \hline
\end{tabular}
\caption{Đặc tả use case UC-004}
\end{table}

\begin{figure}[H]
\centering
\begin{tikzpicture}[>=Stealth, node distance=0.9cm and 1.4cm,
  act/.style={rectangle, rounded corners=4pt, draw, minimum width=3.6cm, minimum height=0.8cm, align=center, font=\small, fill=blue!5},
  dec/.style={diamond, aspect=2.2, draw, align=center, font=\scriptsize, inner sep=1pt, fill=yellow!15},
  st/.style={circle, fill=black, minimum size=0.35cm},
  en/.style={circle, draw, double, minimum size=0.35cm}]
\node[st] (s) {};
\node[act, below=0.5cm of s] (a1) {Mở trang Add new category};
\node[act, below=0.5cm of a1] (a2) {Nhập dữ liệu vào form};
\node[dec, below=0.6cm of a2] (d0) {Bấm\\Cancel?};
\node[act, below=0.7cm of d0] (a3) {Bấm Create category};
\node[dec, below=0.6cm of a3] (d1) {Name, Parent\\hợp lệ?};
\node[dec, below=0.7cm of d1] (d2) {ID number\\không trùng?};
\node[dec, below=0.7cm of d2] (d3) {Name\\$\leq$ 255 ký tự?};
\node[act, below=0.7cm of d3] (a4) {Tạo danh mục, chuyển về\\Manage course categories};
\node[en, below=0.5cm of a4] (e) {};
\node[act, right=2.6cm of d0] (c1) {Quay lại trang quản lý,\\không tạo danh mục};
\node[act, right=2.6cm of d1] (er1) {Hiển thị lỗi\\tại trường};
\node[act, right=2.6cm of d2] (er2) {Hiển thị lỗi\\ID number trùng};
\node[act, right=2.6cm of d3] (er3) {Trang lỗi\\Category name too long};
\draw[->] (s)--(a1); \draw[->] (a1)--(a2); \draw[->] (a2)--(d0);
\draw[->] (d0)--node[right,font=\scriptsize]{Không}(a3);
\draw[->] (d0)--node[above,font=\scriptsize]{Có}(c1);
\draw[->] (a3)--(d1);
\draw[->] (d1)--node[right,font=\scriptsize]{Có}(d2);
\draw[->] (d1)--node[above,font=\scriptsize]{Không}(er1);
\draw[->] (d2)--node[right,font=\scriptsize]{Có}(d3);
\draw[->] (d2)--node[above,font=\scriptsize]{Không}(er2);
\draw[->] (d3)--node[right,font=\scriptsize]{Có}(a4);
\draw[->] (d3)--node[above,font=\scriptsize]{Không}(er3);
\draw[->] (a4)--(e);
\draw[->] (er1.east)--++(0.8,0)|-(a2.east);
\draw[->] (er2.east)--++(1.2,0)|-(a2.east);
\draw[->] (er3.east)--++(1.6,0)|-(a2.east);
\draw[->] (c1.east)--++(2.2,0)|-(e.east);
\end{tikzpicture}
\caption{Activity diagram của UC-004 (các nhánh báo lỗi quay lại bước nhập dữ liệu; nhánh Cancel kết thúc)}
\end{figure}

\subsection{Các biến đầu vào}
\begin{table}[H]
\centering\small
\begin{tabular}{|L{2.6cm}|L{2.0cm}|L{5.2cm}|L{5.2cm}|}
\hline
\rowcolor{gray!15}\textbf{Biến} & \textbf{Bắt buộc?} & \textbf{Ràng buộc} & \textbf{Ghi chú} \\ \hline
Parent category & Có & Chọn từ danh sách (Top, các danh mục có sẵn); không chọn (No selection) thì báo lỗi & Lỗi: \textit{You must supply a value here.} \\ \hline
Category name & Có & 1--255 ký tự; được phép trùng tên; chuỗi chỉ gồm khoảng trắng vẫn được gửi đi (TC-004-003) & Rỗng: \textit{Required}; $>$255 ký tự: trang lỗi \textit{Category name too long} \\ \hline
Category ID number & Không & 0--100 ký tự; không được trùng với danh mục khác & Ô chặn nhập quá 100 ký tự (\texttt{maxlength = 100}); trùng: \textit{ID number is already used for another category} \\ \hline
Description & Không & Văn bản tự do (trình soạn thảo TinyMCE), không có giới hạn & Không có giá trị không hợp lệ qua giao diện; Katalon chỉ kiểm thử trường hợp để trống (xem mục Katalon) \\ \hline
\end{tabular}
\caption{Các biến đầu vào của Feature 004}
\end{table}

\subsection{Kỹ thuật 1: Phân lớp tương đương (ECP)}
\textbf{Cách áp dụng:} với mỗi biến, chia miền giá trị thành lớp hợp lệ (V) và các lớp không hợp lệ (I).
Test case ``tất cả hợp lệ'' (TC-004-001) phủ các lớp hợp lệ của Parent, Name, ID number và Description; mỗi lớp không hợp lệ được kiểm thử
trong một test case riêng, các trường khác giữ giá trị hợp lệ để cô lập nguyên nhân lỗi.

\begin{table}[H]
\centering\footnotesize
\begin{tabular}{|l|L{5.4cm}|L{7.0cm}|}
\hline
\rowcolor{gray!15}\textbf{Biến} & \textbf{Lớp hợp lệ} & \textbf{Lớp không hợp lệ (test case)} \\ \hline
Parent & V1: chọn Top hoặc một danh mục & I1: No selection (007) \\ \hline
Name & V2: 1--255 ký tự (kể cả chuỗi chỉ có khoảng trắng, xem 003) & I2: rỗng (002); I3: $>$255 ký tự (004) \\ \hline
ID number & V3: rỗng; V4: 1--100 ký tự, chưa tồn tại & I4: $>$100 ký tự (005); I5: đã tồn tại (006) \\ \hline
Description & V5: rỗng (Katalon); có nội dung (chỉ thử thủ công) & Không có \\ \hline
\end{tabular}
\caption{Lớp tương đương của Feature 004 (ID rút gọn: 002 nghĩa là TC-004-002)}
\end{table}

\subsection{Kỹ thuật 2: Phân tích giá trị biên (BVA)}
\textbf{Cách áp dụng:} với mỗi biến có miền độ dài liên tục, chọn các giá trị tại và quanh biên: min$-1$, min, min$+1$, max$-1$, max, max$+1$.
Giá trị \emph{không hợp lệ} (min$-1$, max$+1$) trùng với test case ECP nên được tham chiếu lại; các giá trị hợp lệ tạo test case mới (008--017).
Parent là danh sách rời rạc nên không áp dụng BVA; Description là văn bản tự do không có giới hạn trên nên chỉ chọn biên dưới là 0 ký tự (để trống).

\begin{table}[H]
\centering\footnotesize
\begin{tabular}{|l|c|c|c|c|c|c|}
\hline
\rowcolor{gray!15}\textbf{Biến (độ dài)} & \textbf{min$-1$} & \textbf{min} & \textbf{min$+1$} & \textbf{max$-1$} & \textbf{max} & \textbf{max$+1$} \\ \hline
Name (1--255) & 0 (002) & 1 (008) & 2 (009) & 254 (010) & 255 (011) & 256 (004) \\ \hline
ID number (0--100) & -- & 0 (012) & 1 (013) & 99 (015) & 100 (016) & 101 (005) \\ \hline
Description (0--$\infty$) & -- & 0 (017) & -- & -- & -- & -- \\ \hline
\end{tabular}
\caption{Giá trị biên và test case của Feature 004 (số trong ngoặc là TC-004-YYY; ID number có thêm giá trị 2 ký tự ở TC-004-014)}
\end{table}

\subsection{Kỹ thuật 3: Bảng quyết định (Decision table)}
\textbf{Điều kiện:} C1 -- Name hợp lệ (1--255 ký tự); C2 -- Parent đã chọn; C3 -- ID number (E: rỗng, U: có và duy nhất, D: trùng);
C4 -- Parent là danh mục khác Top. \textbf{Hành động:} A1 -- tạo danh mục ở cấp Top; A2 -- tạo danh mục dưới một danh mục khác; A3 -- báo lỗi.
R7 kiểm tra việc Moodle cho phép hai danh mục cùng tên, kể cả cùng parent. Name $>$255 ký tự không đưa vào bảng vì kết quả là trang lỗi riêng,
đã được kiểm thử ở TC-004-004.

\begin{table}[H]
\centering\footnotesize
\begin{tabular}{|l|c|c|c|c|c|c|c|c|}
\hline
\rowcolor{gray!15}\textbf{Luật} & R1 & R2 & R3 & R4 & R5 & R6 & R7 & R8 \\ \hline
C1 Name hợp lệ & Y & Y & Y & N & Y & N & Y & Y \\ \hline
C2 Parent đã chọn & Y & Y & Y & Y & N & N & Y & Y \\ \hline
C3 ID number & E & U & D & -- & -- & -- & E & U \\ \hline
C4 Parent khác Top & N & N & Y & -- & -- & -- & Y & Y \\ \hline
A1 Tạo ở cấp Top & X & X & & & & & & \\ \hline
A2 Tạo dưới danh mục khác & & & & & & & X & X \\ \hline
A3 Báo lỗi & & & X & X & X & X & & \\ \hline
\textbf{Test case} & 012 & 001 & 006 & 002 & 007 & 018 & 019 & 020 \\ \hline
\end{tabular}
\caption{Decision table của Feature 004 (R4 dùng test case ECP 002; R3 và R7 dùng parent là tc004\_root)}
\end{table}
\textit{Ghi chú:} R7 tạo hai danh mục cùng tên \texttt{tc004\_same} dưới \texttt{tc004\_root} và cả hai đều được tạo (hệ thống phân biệt bằng id).

\subsection{Kỹ thuật 4: Use-case testing}
Từ activity diagram, các đường đi (scenario) được liệt kê như sau:

\begin{table}[H]
\centering\footnotesize
\begin{tabular}{|c|L{7.4cm}|L{3.2cm}|c|}
\hline
\rowcolor{gray!15}\textbf{Path} & \textbf{Đường đi trong diagram} & \textbf{Loại} & \textbf{Test case} \\ \hline
P1 & Mở form $\rightarrow$ nhập hợp lệ $\rightarrow$ Create category $\rightarrow$ tạo thành công & Luồng chính & 021 \\ \hline
P2 & Gửi form thiếu Name $\rightarrow$ báo lỗi $\rightarrow$ nhập lại $\rightarrow$ thành công & Luồng thay thế A1 & 022 \\ \hline
P3 & Gửi ID number trùng $\rightarrow$ báo lỗi $\rightarrow$ đổi ID $\rightarrow$ thành công & Luồng thay thế A2 & 023 \\ \hline
P4 & Nhập dữ liệu $\rightarrow$ Cancel $\rightarrow$ không tạo danh mục & Luồng thay thế A5 & 024 \\ \hline
P5 & Tạo danh mục cấp 1 $\rightarrow$ tạo danh mục con của nó & Luồng thay thế A4 & 025 \\ \hline
P6 & Để Parent No selection $\rightarrow$ báo lỗi $\rightarrow$ chọn Top $\rightarrow$ thành công & Luồng thay thế A1 & 026 \\ \hline
\end{tabular}
\caption{Scenario và test case của use-case testing (Feature 004; nhánh A3 -- Name $>$255 ký tự -- được kiểm thử ở TC-004-004)}
\end{table}

\subsection{Danh sách test case của Feature 004}
Tổng cộng \textbf{26 test case}: ECP 001--007, BVA 008--017, Decision table 018--020, Use-case 021--026.
Các test case ghi ``ECP/BVA'' vừa là lớp không hợp lệ vừa là giá trị biên.
Kết quả mong đợi chung: với test case hợp lệ -- chuyển về trang Manage course categories and courses và danh mục mới có trong cây danh mục;
với test case không hợp lệ -- form báo lỗi tại trường tương ứng và danh mục không được tạo.
Giới hạn 100 ký tự của ID number được kiểm tra bằng thuộc tính \texttt{maxlength} của ô nhập; giới hạn 255 ký tự của Name được kiểm tra
bằng cách gửi form và quan sát trang lỗi. Các test case chạy với Parent đã chọn trong ô Parent category (trừ 007, 018, 026 kiểm tra Parent chưa chọn).

{\footnotesize
\begin{longtable}{|l|L{4.4cm}|c|L{6.4cm}|}
\hline
\rowcolor{gray!15}\textbf{ID} & \textbf{Tên test case} & \textbf{Kỹ thuật} & \textbf{Mô tả} \\ \hline
\endhead
TC-004-001 & Tạo danh mục với dữ liệu hợp lệ & ECP & Parent Top, Name tc004\_\allowbreak\{\}valid, ID number tc004\_\allowbreak\{\}id001; Description để trống. \\ \hline
TC-004-002 & Category name rỗng & ECP & Lớp không hợp lệ I2: để trống Category name (Parent tc004\_\allowbreak\{\}root). \\ \hline
TC-004-003 & Category name chỉ có khoảng trắng & ECP & Lớp V2 đặc biệt: Name chỉ có 3 dấu cách; Moodle không báo Required và chuyển về trang quản lý. \\ \hline
TC-004-004 & Category name 256 ký tự (max+1) & ECP/BVA & Lớp I3: Name dài 256 ký tự; Moodle hiện trang lỗi ``Category name too long''. \\ \hline
TC-004-005 & Category ID number 101 ký tự (max+1) & ECP/BVA & Lớp I4: nhập 101 ký tự; ô chỉ nhận tối đa 100 ký tự (maxlength = 100). \\ \hline
TC-004-006 & Category ID number đã tồn tại & ECP & Lớp I5: ID number tc004\_\allowbreak\{\}dup trùng ID của danh mục tc004\_\allowbreak\{\}base. \\ \hline
TC-004-007 & Parent category = No selection & ECP & Lớp I1: bỏ chọn Parent category. \\ \hline
TC-004-008 & Category name 1 ký tự & BVA & Biên hợp lệ của Name: độ dài 1. \\ \hline
TC-004-009 & Category name 2 ký tự & BVA & Biên hợp lệ của Name: độ dài 2. \\ \hline
TC-004-010 & Category name 254 ký tự & BVA & Biên hợp lệ của Name: độ dài 254. \\ \hline
TC-004-011 & Category name 255 ký tự & BVA & Biên hợp lệ của Name: độ dài 255. \\ \hline
TC-004-012 & Category ID number 0 ký tự (để trống) & BVA & Biên hợp lệ của ID number: độ dài 0; Parent Top. \\ \hline
TC-004-013 & Category ID number 1 ký tự & BVA & Biên hợp lệ của ID number: độ dài 1. \\ \hline
TC-004-014 & Category ID number 2 ký tự & BVA & Biên hợp lệ của ID number: độ dài 2. \\ \hline
TC-004-015 & Category ID number 99 ký tự & BVA & Biên hợp lệ của ID number: độ dài 99. \\ \hline
TC-004-016 & Category ID number 100 ký tự & BVA & Biên hợp lệ của ID number: độ dài 100. \\ \hline
TC-004-017 & Description để trống (0 ký tự) & BVA & Biên dưới của Description: không nhập Description, danh mục vẫn tạo được. \\ \hline
TC-004-018 & R6: Name rỗng và Parent No selection & DT & R6: Name rỗng và Parent chưa chọn; hiện cả hai lỗi. \\ \hline
TC-004-019 & R7: Hai danh mục trùng tên & DT & R7: tạo hai danh mục cùng tên, cùng parent; cả hai đều tạo được. \\ \hline
TC-004-020 & R8: Parent là danh mục khác Top & DT & R8: Parent là tc004\_\allowbreak\{\}root; kiểm tra quan hệ cha-con. \\ \hline
TC-004-021 & UC-P1: Luồng chính tạo danh mục & UC & Từ form, chọn Top, nhập hợp lệ, tạo thành công. \\ \hline
TC-004-022 & UC-P2: Bỏ trống Name rồi nhập lại & UC & Gửi form rỗng Name, nhận lỗi, nhập Name rồi gửi lại. \\ \hline
TC-004-023 & UC-P3: ID number trùng rồi đổi & UC & Dùng ID number trùng, nhận lỗi, đổi ID rồi gửi lại. \\ \hline
TC-004-024 & UC-P4: Hủy tạo danh mục (Cancel) & UC & Nhập Name rồi bấm Cancel; không tạo danh mục. \\ \hline
TC-004-025 & UC-P5: Tạo danh mục con của danh mục con & UC & Tạo danh mục cấp 1 rồi tạo danh mục con của nó (cấp 2). \\ \hline
TC-004-026 & UC-P6: Parent No selection, báo lỗi, chọn Top & UC & Bỏ chọn Parent, nhận lỗi, chọn Top rồi gửi lại. \\ \hline
\end{longtable}
}


\subsection{Test case mẫu (đầy đủ bước và kết quả mong đợi)}
Năm test case dưới đây minh họa cách ghi trong file \texttt{.xlsx}; các test case còn lại viết theo cùng khuôn.
Dữ liệu dùng tiền tố \texttt{tc004\_} để dễ dọn dẹp; tên và ID number là giá trị cố định, đúng như trong script Katalon.

\bigskip\noindent\textbf{TC-004-001 -- Tạo danh mục với dữ liệu hợp lệ}\\
\noindent\textit{Precondition:} Moodle demo (Docker) đang chạy tại http:/\allowbreak{}/\allowbreak{}127.0.0.1:9501; giao diện Moodle và trình duyệt là English; trình duyệt đã đăng nhập bằng admin (admin/\allowbreak{}admin, mật khẩu không bị thay đổi); đã tạo thủ công danh mục cấp Top tên tc004\_\allowbreak\{\}root, bên trong có danh mục con tc004\_\allowbreak\{\}base với Category ID number là tc004\_\allowbreak\{\}dup; chưa có danh mục nào khác dùng tên hoặc ID number của test này (xem mục dọn dữ liệu).
\smallskip

{\footnotesize
\begin{longtable}{|c|L{3.0cm}|L{5.1cm}|L{5.6cm}|}
\hline
\rowcolor{gray!15}\textbf{\#} & \textbf{Bước} & \textbf{Mô tả} & \textbf{Kết quả mong đợi} \\ \hline
\endhead
1 & Mở form Add new category & Mở \textless{}BASE\_\allowbreak\{\}URL\textgreater{}/\allowbreak{}course/\allowbreak{}editcategory.php?parent=0. & Form Add new category hiển thị. \\ \hline
2 & Chọn Parent category & Trong ô Parent category chọn Top. & Ô Parent category hiển thị Top. \\ \hline
3 & Nhập Category name & Nhập tc004\_\allowbreak\{\}valid vào ô Category name. & Ô hiển thị đúng giá trị đã nhập. \\ \hline
4 & Nhập Category ID number & Nhập tc004\_\allowbreak\{\}id001 vào ô Category ID number. & Ô hiển thị đúng giá trị đã nhập. \\ \hline
5 & Bấm Create category & Bấm nút Create category. & -- \\ \hline
6 & Kiểm tra kết quả & Quan sát trang sau khi gửi. & Chuyển tới trang Manage course categories and courses; danh mục tc004\_\allowbreak\{\}valid xuất hiện trong cây danh mục. \\ \hline
\end{longtable}
}

\bigskip\noindent\textbf{TC-004-003 -- Category name chỉ có khoảng trắng}\\
\noindent\textit{Precondition:} Moodle demo (Docker) đang chạy tại http:/\allowbreak{}/\allowbreak{}127.0.0.1:9501; giao diện Moodle và trình duyệt là English; trình duyệt đã đăng nhập bằng admin (admin/\allowbreak{}admin, mật khẩu không bị thay đổi); đã tạo thủ công danh mục cấp Top tên tc004\_\allowbreak\{\}root, bên trong có danh mục con tc004\_\allowbreak\{\}base với Category ID number là tc004\_\allowbreak\{\}dup; chưa có danh mục nào khác dùng tên hoặc ID number của test này (xem mục dọn dữ liệu).
\smallskip

{\footnotesize
\begin{longtable}{|c|L{3.0cm}|L{5.1cm}|L{5.6cm}|}
\hline
\rowcolor{gray!15}\textbf{\#} & \textbf{Bước} & \textbf{Mô tả} & \textbf{Kết quả mong đợi} \\ \hline
\endhead
1 & Mở form Add new category & Mở \textless{}BASE\_\allowbreak\{\}URL\textgreater{}/\allowbreak{}course/\allowbreak{}editcategory.php?parent=0. & Form Add new category hiển thị. \\ \hline
2 & Chọn Parent category & Trong ô Parent category chọn tc004\_\allowbreak\{\}root. & Ô Parent category hiển thị tc004\_\allowbreak\{\}root. \\ \hline
3 & Nhập Category name & Nhập 3 dấu cách vào ô Category name. & Ô Category name chứa 3 dấu cách. \\ \hline
4 & Bấm Create category & Bấm nút Create category. & -- \\ \hline
5 & Kiểm tra kết quả & Quan sát trang sau khi gửi. & Không hiện lỗi Required; chuyển tới trang Manage course categories and courses. \\ \hline
6 & Mở tc004\_\allowbreak\{\}root & Bấm link tc004\_\allowbreak\{\}root trong cây danh mục. & Trang hiển thị nội dung danh mục tc004\_\allowbreak\{\}root. \\ \hline
7 & Kiểm tra danh mục con & Quan sát danh sách danh mục con. & Danh mục tc004\_\allowbreak\{\}base hiển thị trong tc004\_\allowbreak\{\}root. \\ \hline
\end{longtable}
}

\bigskip\noindent\textbf{TC-004-004 -- Category name 256 ký tự (max+1)}\\
\noindent\textit{Precondition:} Moodle demo (Docker) đang chạy tại http:/\allowbreak{}/\allowbreak{}127.0.0.1:9501; giao diện Moodle và trình duyệt là English; trình duyệt đã đăng nhập bằng admin (admin/\allowbreak{}admin, mật khẩu không bị thay đổi); đã tạo thủ công danh mục cấp Top tên tc004\_\allowbreak\{\}root, bên trong có danh mục con tc004\_\allowbreak\{\}base với Category ID number là tc004\_\allowbreak\{\}dup; chưa có danh mục nào khác dùng tên hoặc ID number của test này (xem mục dọn dữ liệu).
\smallskip

{\footnotesize
\begin{longtable}{|c|L{3.0cm}|L{5.1cm}|L{5.6cm}|}
\hline
\rowcolor{gray!15}\textbf{\#} & \textbf{Bước} & \textbf{Mô tả} & \textbf{Kết quả mong đợi} \\ \hline
\endhead
1 & Mở form Add new category & Mở \textless{}BASE\_\allowbreak\{\}URL\textgreater{}/\allowbreak{}course/\allowbreak{}editcategory.php?parent=0. & Form Add new category hiển thị. \\ \hline
2 & Chọn Parent category & Trong ô Parent category chọn tc004\_\allowbreak\{\}root. & Ô Parent category hiển thị tc004\_\allowbreak\{\}root. \\ \hline
3 & Nhập Category name & Nhập chuỗi 256 ký tự (bắt đầu bằng tc004\_\allowbreak\{\}n256\_\allowbreak\{\}, phần còn lại là chữ a) vào ô Category name. & Ô Category name nhận đủ 256 ký tự (ô không chặn nhập). \\ \hline
4 & Bấm Create category & Bấm nút Create category. & -- \\ \hline
5 & Kiểm tra lỗi & Quan sát trang sau khi gửi. & Không tạo danh mục; Moodle hiển thị trang Error với dòng ``Category name too long''. \\ \hline
\end{longtable}
}

\bigskip\noindent\textbf{TC-004-006 -- Category ID number đã tồn tại}\\
\noindent\textit{Precondition:} Moodle demo (Docker) đang chạy tại http:/\allowbreak{}/\allowbreak{}127.0.0.1:9501; giao diện Moodle và trình duyệt là English; trình duyệt đã đăng nhập bằng admin (admin/\allowbreak{}admin, mật khẩu không bị thay đổi); đã tạo thủ công danh mục cấp Top tên tc004\_\allowbreak\{\}root, bên trong có danh mục con tc004\_\allowbreak\{\}base với Category ID number là tc004\_\allowbreak\{\}dup; chưa có danh mục nào khác dùng tên hoặc ID number của test này (xem mục dọn dữ liệu).
\smallskip

{\footnotesize
\begin{longtable}{|c|L{3.0cm}|L{5.1cm}|L{5.6cm}|}
\hline
\rowcolor{gray!15}\textbf{\#} & \textbf{Bước} & \textbf{Mô tả} & \textbf{Kết quả mong đợi} \\ \hline
\endhead
1 & Mở form Add new category & Mở \textless{}BASE\_\allowbreak\{\}URL\textgreater{}/\allowbreak{}course/\allowbreak{}editcategory.php?parent=0. & Form Add new category hiển thị. \\ \hline
2 & Chọn Parent category & Trong ô Parent category chọn tc004\_\allowbreak\{\}root. & Ô Parent category hiển thị tc004\_\allowbreak\{\}root. \\ \hline
3 & Nhập Category name & Nhập tc004\_\allowbreak\{\}dupid vào ô Category name. & Ô hiển thị đúng giá trị đã nhập. \\ \hline
4 & Nhập Category ID number & Nhập tc004\_\allowbreak\{\}dup (trùng ID của danh mục tc004\_\allowbreak\{\}base). & Ô hiển thị đúng giá trị đã nhập. \\ \hline
5 & Bấm Create category & Bấm nút Create category. & -- \\ \hline
6 & Kiểm tra lỗi & Quan sát form sau khi gửi. & Không tạo danh mục; vẫn ở trang Add new category; hiện lỗi ``ID number is already used for another category''. \\ \hline
\end{longtable}
}

\bigskip\noindent\textbf{TC-004-025 -- UC-P5: Tạo danh mục con của danh mục con}\\
\noindent\textit{Precondition:} Moodle demo (Docker) đang chạy tại http:/\allowbreak{}/\allowbreak{}127.0.0.1:9501; giao diện Moodle và trình duyệt là English; trình duyệt đã đăng nhập bằng admin (admin/\allowbreak{}admin, mật khẩu không bị thay đổi); đã tạo thủ công danh mục cấp Top tên tc004\_\allowbreak\{\}root, bên trong có danh mục con tc004\_\allowbreak\{\}base với Category ID number là tc004\_\allowbreak\{\}dup; chưa có danh mục nào khác dùng tên hoặc ID number của test này (xem mục dọn dữ liệu).
\smallskip

{\footnotesize
\begin{longtable}{|c|L{3.0cm}|L{5.1cm}|L{5.6cm}|}
\hline
\rowcolor{gray!15}\textbf{\#} & \textbf{Bước} & \textbf{Mô tả} & \textbf{Kết quả mong đợi} \\ \hline
\endhead
1 & Mở form Add new category & Mở \textless{}BASE\_\allowbreak\{\}URL\textgreater{}/\allowbreak{}course/\allowbreak{}editcategory.php?parent=0. & Form Add new category hiển thị. \\ \hline
2 & Chọn Parent category & Trong ô Parent category chọn tc004\_\allowbreak\{\}root. & Ô Parent category hiển thị tc004\_\allowbreak\{\}root. \\ \hline
3 & Nhập Category name & Nhập tc004\_\allowbreak\{\}p5a vào ô Category name. & Ô hiển thị đúng giá trị đã nhập. \\ \hline
4 & Bấm Create category & Bấm nút Create category. & -- \\ \hline
5 & Kiểm tra danh mục cấp 1 & Quan sát trang sau khi gửi. & Chuyển tới trang Manage course categories and courses; tc004\_\allowbreak\{\}p5a được tạo dưới tc004\_\allowbreak\{\}root. \\ \hline
6 & Mở form tạo danh mục con & Bấm link Create new category; trong ô Parent category chọn tc004\_\allowbreak\{\}p5a. & Form hiển thị; ô Parent category hiển thị tc004\_\allowbreak\{\}p5a. \\ \hline
7 & Nhập Category name & Nhập tc004\_\allowbreak\{\}p5b vào ô Category name. & Ô hiển thị đúng giá trị đã nhập. \\ \hline
8 & Bấm Create category & Bấm nút Create category. & -- \\ \hline
9 & Kiểm tra danh mục cấp 2 & Bấm tc004\_\allowbreak\{\}p5a trong cây danh mục. & Danh mục tc004\_\allowbreak\{\}p5b nằm dưới tc004\_\allowbreak\{\}p5a. \\ \hline
\end{longtable}
}


\section{Feature 005: ---}
% ================= 7 =================
\section{Triển khai test case bằng Katalon Recorder}

\subsection{Cấu trúc test suite}
Bộ test của Feature 001 nằm trong một file \texttt{F001\_Add\_New\_User.krecorder} (định dạng XHTML của Katalon Recorder)
gồm 51 test case:
\begin{itemize}[noitemsep]
  \item \textbf{TC-000-Setup-Login} (chạy đầu tiên, chỉ chạy một lần): đăng nhập admin bằng tài khoản có sẵn, không đổi mật khẩu.
  \item \textbf{TC-001-001 $\dots$ TC-001-050}: mỗi test case của mục 5.8 là một test case Katalon, tên trùng ID.
        Mỗi test case tự lưu \texttt{BASE\_URL} và biến \texttt{RUN} (\texttt{Date.now().toString(36)}) rồi ghép \texttt{\$\{RUN\}} vào
        username và email, nên chạy lại nhiều lần không bị trùng dữ liệu.
\end{itemize}
Mỗi test case có ít nhất một lệnh \texttt{verify...} (\texttt{verifyLocation}, \texttt{verifyTextPresent},
\texttt{verifyElementPresent}, \texttt{verifyAttribute}, \texttt{verifyChecked}, \dots). Các test case chỉ kiểm tra giới hạn
độ dài (TC-001-008, 010, 012, 016, 017) không gửi form mà đọc thuộc tính \texttt{maxlength} của ô nhập bằng
\texttt{verifyAttribute}. Locator dùng \texttt{id=} của Moodle (\texttt{id\_username}, \texttt{id\_newpassword},
\texttt{id\_submitbutton}, \texttt{id\_error\_username}, \dots).

\subsection{Script Setup (đăng nhập, chỉ chạy một lần đầu suite)}
\begin{lstlisting}
store                 | http://127.0.0.1:9501 | BASE_URL
open                  | ${BASE_URL}/login/index.php
type                  | id=username | admin
type                  | id=password | admin
clickAndWait          | id=loginbtn
verifyElementPresent  | id=user-menu-toggle
\end{lstlisting}

\subsection{TC-001-001 (dữ liệu hợp lệ)}
\begin{lstlisting}
store              | http://127.0.0.1:9501 | BASE_URL
storeEval          | Date.now().toString(36) | RUN
open               | ${BASE_URL}/user/editadvanced.php?id=-1
type               | id=id_username | tc001_valid${RUN}
click              | xpath=//div[@data-passwordunmaskid='id_newpassword']//a[@data-passwordunmask='edit']
type               | id=id_newpassword | Test@1234
type               | id=id_firstname | Nguyen
type               | id=id_lastname | Van A
type               | id=id_email | tc001_valid${RUN}@example.com
type               | id=id_city | Ho Chi Minh
clickAndWait       | id=id_submitbutton
verifyLocation     | glob:*/admin/user.php*
verifyTextPresent  | Changes saved
verifyTextPresent  | tc001_valid${RUN}@example.com
\end{lstlisting}

\subsection{Feature 002: Add a new course}
Bộ test của Feature 002 nằm trong file \texttt{F002-Add a new course.krecorder} (test suite \texttt{TC-002}), gồm
\textbf{38 test case TC-002-001 $\dots$ TC-002-038}, tên trùng ID trong file \texttt{.xlsx}. Các test case được ghi bằng chức năng
\textit{Record} của Katalon Recorder rồi bổ sung lệnh kiểm tra.
\begin{itemize}[noitemsep]
  \item \textbf{Điều kiện trước khi chạy:} admin đã đăng nhập sẵn trong trình duyệt (bước chuẩn bị, không phải đối tượng kiểm thử),
        Moodle chạy tại \texttt{http://127.0.0.1:9501}.
  \item \textbf{Chuẩn bị dữ liệu:} course có short name \texttt{EXIST01} (dùng cho TC-002-008, 032, 037) được tạo ngay trong phần
        đầu của TC-002-008; course có ID number \texttt{ID-EXIST} được tạo trong phần đầu của TC-002-013.
  \item \textbf{Kiểm tra kết quả:} mỗi test case có ít nhất một lệnh \texttt{verifyText}:
        \begin{itemize}[noitemsep]
          \item Test case hợp lệ: mở \textit{My courses} và kiểm tra full name (hoặc category) của course vừa tạo.
          \item Test case lỗi: kiểm tra nội dung thông báo tại \texttt{id\_error\_fullname}, \texttt{id\_error\_shortname},
                \texttt{id\_error\_category}, \texttt{id\_error\_enddate} hoặc \texttt{id\_error\_idnumber}.
          \item TC-002-031 và 032 (decision table) kiểm tra đồng thời 3 và 2 thông báo lỗi.
        \end{itemize}
\end{itemize}

\begin{table}[H]
\centering\footnotesize
\begin{tabular}{|L{3.2cm}|L{4.8cm}|L{6.2cm}|}
\hline
\rowcolor{gray!15}\textbf{Nhóm test case} & \textbf{Locator được verify} & \textbf{Giá trị mong đợi} \\ \hline
001, 002, 004, 005, 007, 012, 014--025, 027--029, 033, 036--038 & Thẻ course trong \textit{My courses} & Full name của course (TC-002-004: ``Category 1'') \\ \hline
003, 031 & \texttt{id=id\_error\_fullname} & Missing full name \\ \hline
006, 031 & \texttt{id=id\_error\_shortname} & Missing short name \\ \hline
008, 032 & \texttt{id=id\_error\_shortname} & Short name is already used for another course (\dots) \\ \hline
009, 031 & \texttt{id=id\_error\_category} & You must supply a value here. \\ \hline
010, 011, 026, 030, 032 & \texttt{id=id\_error\_enddate} & The course end date must be after the start date. \\ \hline
013 & \texttt{id=id\_error\_idnumber} & ID number is already used for another course (\dots) \\ \hline
034 & Danh sách course của trang quản lý & TC002 UC Return \\ \hline
035 & Vùng danh sách course & No courses (course không được tạo sau Cancel) \\ \hline
\end{tabular}
\caption{Lệnh verify của 38 test case Feature 002}
\end{table}

\textbf{TC-002-001 (trích script đã ghi, bỏ các lệnh click trung gian):}
\begin{lstlisting}
open           | http://127.0.0.1:9501/course/edit.php?category=0 |
type           | id=id_fullname    | TC002 Course 001
type           | id=id_shortname   | TC002-001
type           | id=form_autocomplete_input-... | Ca
click          | id=form_autocomplete_suggestions-...-0 |
select         | id=id_category    | label=Category 1
select         | id=id_startdate_day   | label=1
select         | id=id_startdate_month | label=January
select         | id=id_startdate_hour  | label=08
select         | id=id_enddate_day     | label=1
select         | id=id_enddate_month   | label=January
type           | id=id_idnumber    | ID-002-001
click          | id=id_saveanddisplay |
open           | http://127.0.0.1:9501/my/courses.php |
verifyText     | xpath=//div[@id='course-info-container-7-4']/div/div/a/span[3]/span[2] | TC002 Course 001
\end{lstlisting}

\textbf{TC-002-008 (short name trùng, phần đầu tạo course EXIST01):}
\begin{lstlisting}
open           | http://127.0.0.1:9501/course/edit.php?category=0 |
type           | id=id_fullname    | 1
type           | id=id_shortname   | EXIST01
click          | id=id_saveanddisplay |
open           | http://127.0.0.1:9501/course/edit.php?category=0 |
type           | id=id_fullname    | TC002 Course 008
type           | id=id_shortname   | EXIST01
select         | (start/end date: xem TC-002-001) |
click          | id=id_saveanddisplay |
verifyText     | id=id_error_shortname | Short name is already used for another course (1)
\end{lstlisting}

\textbf{TC-002-004 (full name chỉ có khoảng trắng: kết quả thực tế là course được tạo):}
\begin{lstlisting}
open           | http://127.0.0.1:9501/course/edit.php?category=0 |
type           | id=id_fullname    | (3 khoang trang)
type           | id=id_shortname   | TC002-004
select         | (start/end date: xem TC-002-001) |
click          | id=id_saveanddisplay |
open           | http://127.0.0.1:9501/my/courses.php |
verifyText     | xpath=//div[@id='course-info-container-10-4']/div/div/div/span[2] | Category 1
\end{lstlisting}

\textbf{Dọn dữ liệu:} sau khi chạy, vào \textit{Site administration $>$ Courses $>$ Manage courses and categories $>$ Category 1},
xóa các course do test tạo (short name bắt đầu bằng \texttt{TC002}, \texttt{EXIST01}, \texttt{ID-EXIST} và các course BVA
có short name/ID number 1--2 ký tự). Phải dọn trước khi chạy lại suite vì short name cố định sẽ bị trùng.
\textbf{Không xóa course hay category có sẵn của demo.}

\subsection{TC-003-001}
\subsection{Feature 004: TC-004-001 và các test case tiêu biểu}
Bộ test của Feature 004 gồm \textbf{26 test case} Katalon \texttt{TC-004-001} đến \texttt{TC-004-026} trong file
\texttt{F004\_Add\_new\_catogory.krecorder} (suite \texttt{F004\_TestSuite\_Optimized-v6}). Suite \emph{không có} script Setup/Cleanup:
dữ liệu nền (\texttt{tc004\_root}, \texttt{tc004\_base}) được tạo thủ công theo Precondition và admin phải đăng nhập sẵn trong trình duyệt.
Mỗi test case độc lập, tên và ID number được ghi cố định; các biến Katalon (ví dụ \texttt{LEN256}, \texttt{ID101}, \texttt{P5AID}) chỉ dùng trong phạm vi
một test case. Mỗi test case có ít nhất một lệnh \texttt{verify...} (\texttt{verifyTextPresent}, \texttt{verifyText}, \texttt{verifyElementPresent},
\texttt{verifyElementNotPresent}, \texttt{verifyAttribute}). \textbf{Giới hạn:} Description là trình soạn thảo TinyMCE (iframe) nên script không nhập
Description; TC-004-001 và TC-004-017 chỉ kiểm thử trường hợp để trống, trường hợp có nội dung chưa được kiểm thử tự động.

\textbf{TC-004-001 (dữ liệu hợp lệ):}
\begin{lstlisting}
open               | http://127.0.0.1:9501/course/editcategory.php?parent=0
select             | id=id_parent | label=Top
type               | id=id_name | tc004_valid
type               | id=id_idnumber | tc004_id001
clickAndWait       | id=id_submitbutton
verifyTextPresent  | tc004_valid
\end{lstlisting}

\textbf{TC-004-003 (Name chỉ có khoảng trắng; giá trị của lệnh \texttt{type} là 3 dấu cách):}
\begin{lstlisting}
open               | http://127.0.0.1:9501/course/editcategory.php?parent=0
select             | id=id_parent | label=tc004_root
type               | id=id_name |    
clickAndWait       | id=id_submitbutton
clickAndWait       | link=tc004_root
verifyTextPresent  | tc004_base
\end{lstlisting}

\textbf{TC-004-025 (danh mục con của danh mục con):}
\begin{lstlisting}
open                | http://127.0.0.1:9501/course/editcategory.php?parent=0
select              | id=id_parent | label=tc004_root
type                | id=id_name | tc004_p5a
clickAndWait        | id=id_submitbutton
verifyTextPresent   | tc004_p5a
storeAttribute      | link=tc004_p5a@href | P5AHREF
storeEval           | '${P5AHREF}'.split('categoryid=')[1].split('&')[0] | P5AID
clickAndWait        | link=Create new category
select              | id=id_parent | value=${P5AID}
type                | id=id_name | tc004_p5b
clickAndWait        | id=id_submitbutton
clickAndWait        | link=tc004_p5a
verifyElementPresent | link=tc004_p5b
\end{lstlisting}

\textbf{Dọn dữ liệu Feature 004:} các danh mục do test tạo có tên bắt đầu bằng \texttt{tc004\_}. Trước khi chạy lại:
(1) xóa bốn danh mục cấp Top \texttt{tc004\_valid} (TC-004-001), \texttt{tc004\_id0} (TC-004-012), \texttt{tc004\_p1} (TC-004-021),
\texttt{tc004\_p6} (TC-004-026); (2) xóa \texttt{tc004\_root} cùng toàn bộ danh mục con (\textit{Delete $>$ Delete all}), gồm cả danh mục tên rỗng do
TC-004-003 tạo; (3) tạo lại \texttt{tc004\_root} và \texttt{tc004\_base} (ID number \texttt{tc004\_dup}). Các ID number cố định cần được giải phóng:
\texttt{tc004\_id001}, \texttt{q}, \texttt{qq}, hai ID dài 99 và 100 ký tự, \texttt{tc004\_idchild}, \texttt{tc004\_p1id}, \texttt{tc004\_p3id}.
\subsection{TC-005-001}




\subsection{Dọn dữ liệu sau khi chạy}
Các user do test tạo có username/email chứa \texttt{tc001\_}. Để dọn: vào Browse list of users, lọc theo email chứa
\texttt{tc001\_}, chọn tất cả rồi dùng \textit{Bulk user actions $>$ Delete}; \textbf{không chạm vào user mặc định} của demo.
Ba test case dùng giá trị cố định (TC-001-018, TC-001-019, TC-001-034) phải được dọn trước khi chạy lại.


\subsection{Chạy test suite, tải file và log}
\begin{enumerate}[noitemsep]
  \item Import \texttt{F001\_Add\_New\_User.krecorder} vào Katalon Recorder, chạy toàn bộ suite (Setup trước) trên
        \textbf{ít nhất một trình duyệt} (ví dụ Chrome) \textbf{trước khi tải về}; sửa các test case lỗi do locator hoặc dữ liệu.
  \item Tải test suite ở định dạng \texttt{.krecorder} (Export/Save). \textbf{Log chạy được tải kèm test suite}, vì vậy phải chạy xong rồi mới tải.
  \item \textbf{Chạy lại trên một máy khác} để chắc chắn test suite chạy ổn định, rồi mới nộp.
  \item Mỗi test case có ít nhất một lệnh \texttt{verify...}.
\end{enumerate}
\textbf{Phân biệt hai loại mô tả:} file \texttt{.xlsx} là kịch bản kiểm thử \emph{thủ công} (từng bước bằng lời, có kết quả mong đợi),
còn các bảng lệnh ở trên chỉ là script Katalon tương ứng để tự động hóa.

% ================= 8 =================

% ================= 9 =================
\section{AI Usage Declaration}

\subsection{[Nguyễn Đức Phương]}
\textbf{Công cụ:} Claude (Anthropic), model Claude Sonnet 5.5, dùng qua ứng dụng Claude (bản desktop).

\noindent\textbf{Prompt đã sử dụng} (nguyên văn):
\begin{enumerate}[noitemsep]
  \item ``hãy phân tích đề bài này chi tiết cho tôi''
  \item ``nhóm tôi có 5 người, hãy viết báo cáo cho tôi và cho ví dụ 1 vài test case và cách thực hiện''
\end{enumerate}

\noindent\textbf{Kết quả, nhận xét và đề xuất của AI:}
\begin{itemize}[noitemsep]
  \item Phân tích đề bài, thang điểm và các ràng buộc.
  \item Soạn khung báo cáo, phân tích ECP, BVA, decision table, use-case cho Feature 001, 50 test case (6 test case chi tiết),
        script Katalon mẫu và file \texttt{.xlsx} mẫu.
  \item Các giới hạn độ dài, locator và thông báo lỗi do AI đề xuất cần được nhóm đối chiếu với hệ thống Moodle thật trước khi dùng.
\end{itemize}

\subsection{[Lê Sĩ Nhật Quang]}
\textbf{Công cụ:} Claude (Anthropic), model Claude Opus 5.5, dùng qua ứng dụng Claude (bản desktop, tab Code).

\noindent\textbf{Prompt đã sử dụng} (nguyên văn):
\begin{enumerate}[noitemsep]
  \item ``Hướng dẫn tôi thực hiện bài tập lớn 2 dựa trên file mô tả, template file excel'' (kèm file mô tả, file Excel mẫu)
  \item ``chỉnh sửa lại báo cáo với phần cập nhật sau khi đã chạy katalon recorder của tôi với phần chức năng add a new course'' (kèm file 	exttt{.xlsx} đã cập nhật)
\end{enumerate}

\noindent\textbf{Kết quả, nhận xét và đề xuất của AI:}
\begin{itemize}[noitemsep]
  \item Hướng dẫn phân tích Feature 002, sinh 38 test case trong sheet \textit{Add a new course} của file \texttt{.xlsx}.
  \item Soạn mục Feature 002 của báo cáo (use case, activity diagram, ECP, BVA, decision table, use-case testing) và script Katalon mẫu.
  \item Giới hạn độ dài, thông báo lỗi và locator do AI đề xuất ban đầu; thành viên đã tự ghi và chạy script Katalon trên Moodle demo, cập nhật kết quả thực tế vào file 	exttt{.xlsx}, sau đó AI cập nhật báo cáo theo kết quả này.
\end{itemize}
\subsection{[Thành viên 3]}
\todo{Điền công cụ AI và prompt, hoặc câu xác nhận không dùng AI.}
\subsection{[Đào Xuân Sơn]}
\textbf{Công cụ:} Claude (Anthropic), model Claude Sonnet 5.5, dùng qua giao diện chat của ứng dụng Claude.

\noindent\textbf{Prompt đã sử dụng} (nguyên văn, kèm các file đề bài và báo cáo nhóm):
\begin{enumerate}[noitemsep]
  \item ``hướng dẫn chi tiết làm''
  \item ``Cho tôi vài testcase''
  \item ``sửa lại báo cáo cho đồng bộ''
  \item ``tc đúng rồi hãy làm báo cáo đồng bộ với testcase''
  \item ``hãy thêm vào report''
\end{enumerate}

\noindent\textbf{Nội dung AI hỗ trợ:}
\begin{itemize}[noitemsep]
  \item Hướng dẫn các bước làm Project 2 dựa trên đề bài.
  \item Gợi ý thiết kế test case cho Feature 004 theo bốn kỹ thuật (ECP, BVA, decision table, use-case testing); ban đầu cho feature Add an assignment, sau đó cho Add new category.
  \item Soạn và định dạng LaTeX phần Feature 004 trong báo cáo này, đối chiếu với file \texttt{F004\_Add\_new\_catogory.krecorder} của nhóm.
\end{itemize}
\subsection{[Thành viên 5]}
\todo{Điền công cụ AI và prompt, hoặc câu xác nhận không dùng AI.}

\end{document}
